\documentclass[13pt,a4paper]{extarticle}

\usepackage[T5]{fontenc}
\usepackage[utf8]{inputenc}
\usepackage[vietnamese]{babel}
\usepackage{geometry}
\usepackage{graphicx}
\usepackage{mathptmx}
\usepackage{indentfirst}
\usepackage{tocloft}
\usepackage{hyphenat}
\usepackage{hyperref}
\usepackage{float}
\usepackage{amssymb}
\usepackage{amsmath}
\usepackage{xcolor}
\usepackage{array}
\usepackage{tabularx}
\usepackage{longtable}
\usepackage{booktabs}
\usepackage{enumitem}
\usepackage{tikz}
\usepackage{scrextend}
\usepackage{listings}
\usepackage{xurl}
\usepackage{textcomp}
\usepackage{caption}
\usepackage{etoolbox}
\usetikzlibrary{calc,positioning,arrows.meta,shapes.geometric,fit,backgrounds}

\definecolor{myblue}{RGB}{0,72,160}
\definecolor{lightblue}{RGB}{232,242,252}
\definecolor{lightgreen}{RGB}{232,247,238}
\definecolor{linegray}{RGB}{120,120,120}
\newcolumntype{Y}{>{\raggedright\arraybackslash}X}

\hypersetup{
	hidelinks,
	pdftitle={Bao cao do an AgriSmart},
	pdfauthor={Dieu Chinh Sim}
}

\geometry{top=2.5cm,bottom=2.5cm,left=3.5cm,right=2cm,marginparwidth=1.75cm}
\linespread{1.5}
\setlength{\parindent}{1cm}
\setlength{\parskip}{0pt}
\setlength{\emergencystretch}{3em}
\setlist[itemize]{label=-,leftmargin=1.5cm}
\setlist[enumerate]{leftmargin=1.5cm}
\lstset{
	basicstyle=\ttfamily,
	breaklines=true,
	columns=fullflexible,
	keepspaces=true,
	showstringspaces=false,
	xleftmargin=0.7cm
}

\setlength{\cftsubsecindent}{1.5em}
\setlength{\cftsubsubsecindent}{3.5em}
\renewcommand{\cftsubsecdotsep}{\cftdotsep}
\renewcommand{\cftsubsubsecdotsep}{\cftdotsep}
\renewcommand{\cftsecfont}{\bfseries}
\renewcommand{\cftsubsecfont}{\normalfont}
\renewcommand{\cftsubsubsecfont}{\normalfont}
\renewcommand{\cfttoctitlefont}{\fontsize{14pt}{20pt}\selectfont\bfseries}
\renewcommand{\cftaftertoctitle}{}
\renewcommand{\contentsname}{\texorpdfstring{\makebox[\textwidth][c]{MỤC LỤC}}{MỤC LỤC}}
\renewcommand{\listfigurename}{DANH MỤC HÌNH}
\renewcommand{\listtablename}{DANH MỤC BẢNG}
\renewcommand{\refname}{\texorpdfstring{\makebox[\textwidth][c]{DANH MỤC TÀI LIỆU THAM KHẢO}}{DANH MỤC TÀI LIỆU THAM KHẢO}}
\addto\captionsvietnamese{%
	\renewcommand{\contentsname}{\texorpdfstring{\makebox[\textwidth][c]{MỤC LỤC}}{MỤC LỤC}}%
	\renewcommand{\listfigurename}{DANH MỤC HÌNH}%
	\renewcommand{\listtablename}{DANH MỤC BẢNG}%
	\renewcommand{\refname}{\texorpdfstring{\makebox[\textwidth][c]{DANH MỤC TÀI LIỆU THAM KHẢO}}{DANH MỤC TÀI LIỆU THAM KHẢO}}%
	\renewcommand{\figurename}{Hình}%
	\renewcommand{\tablename}{Bảng}%
	\renewcommand{\abstractname}{TÓM TẮT}%
}
\captionsetup{font=normalsize,labelsep=period}

\renewcommand{\cftloftitlefont}{\hfill\fontsize{14pt}{20pt}\selectfont\bfseries}
\renewcommand{\cftafterloftitle}{\hfill}
\renewcommand{\cftlottitlefont}{\hfill\fontsize{14pt}{20pt}\selectfont\bfseries}
\renewcommand{\cftafterlottitle}{\hfill}
\renewcommand{\cftfigpresnum}{Hình~}
\renewcommand{\cfttabpresnum}{Bảng~}
\renewcommand{\cftfigaftersnum}{\quad}
\renewcommand{\cfttabaftersnum}{\quad}
\setlength{\cftfignumwidth}{2.4cm}
\setlength{\cfttabnumwidth}{2.4cm}
\renewcommand{\cftfigdotsep}{\cftdotsep}
\renewcommand{\cfttabdotsep}{\cftdotsep}

\renewcommand{\thesection}{\arabic{section}}
\renewcommand{\thesubsection}{\arabic{section}.\arabic{subsection}}
\renewcommand{\thesubsubsection}{\thesubsection.\arabic{subsubsection}}
\numberwithin{figure}{section}
\numberwithin{table}{section}
\setcounter{secnumdepth}{3}
\setcounter{tocdepth}{3}

\newcommand{\studentname}{Điêu Chính Sim}
\newcommand{\advisorname}{ThS. Nguyễn Văn Hải}

\newcommand{\missinggraphic}[1]{%
	\fbox{%
		\begin{minipage}[c][4.5cm][c]{0.78\linewidth}
			\centering
			\textbf{Thiếu tệp hình}\\
			\texttt{#1}
		\end{minipage}%
	}%
}

\newcommand{\safeincludegraphics}[2][]{%
	\IfFileExists{#2}{\includegraphics[#1]{#2}}{%
		\IfFileExists{#2.pdf}{\includegraphics[#1]{#2.pdf}}{%
			\IfFileExists{#2.png}{\includegraphics[#1]{#2.png}}{%
				\IfFileExists{#2.jpg}{\includegraphics[#1]{#2.jpg}}{%
					\IfFileExists{#2.jpeg}{\includegraphics[#1]{#2.jpeg}}{\missinggraphic{#2}}%
				}%
			}%
		}%
	}%
}

\newif\iffirstchapter
\firstchaptertrue
\newcommand{\chaptersection}[1]{%
	\clearpage
	\iffirstchapter
	\pagenumbering{arabic}
	\setcounter{page}{1}
	\firstchapterfalse
	\fi
	\refstepcounter{section}
	\setcounter{subsection}{0}
	\setcounter{subsubsection}{0}
	\begin{center}
		\fontsize{16pt}{24pt}\selectfont
		\textbf{CHƯƠNG \arabic{section}: #1}
	\end{center}
	\phantomsection
	\addcontentsline{toc}{section}{CHƯƠNG \arabic{section}: #1}
}


\begin{document}
	
	\pagenumbering{roman}
	\setcounter{page}{1}
	\pagestyle{plain}
	\begin{center}
		\fontsize{14pt}{20pt}\selectfont
		\textbf{LỜI CẢM ƠN}
	\end{center}
	
	\vspace{20pt}
	
	Trong suốt quá trình nghiên cứu, phân tích yêu cầu, hiện thực hóa chức năng và hoàn thiện báo cáo cho đồ án \textbf{thiết kế và xây dựng hệ thống giám sát nông nghiệp thông minh}, em đã nhận được nhiều sự hỗ trợ quý báu từ thầy cô, bạn bè và các tài liệu chuyên môn liên quan đến Internet vạn vật, nông nghiệp chính xác và phát triển phần mềm web hiện đại.
	
	Em xin bày tỏ lòng biết ơn sâu sắc tới \advisorname, người đã định hướng phương pháp tiếp cận, góp ý cho em trong quá trình xác định phạm vi, thiết kế kiến trúc và hoàn thiện báo cáo. Những nhận xét sát thực đã giúp em nhìn rõ hơn cách chuyển hóa một bài toán kỹ thuật thành một sản phẩm có giá trị ứng dụng trong thực tiễn sản xuất nông nghiệp.
	
	Em xin chân thành cảm ơn Ban Giám hiệu \textbf{Trường Đại học Tây Bắc}, cùng tập thể quý thầy cô thuộc \textbf{Khoa Khoa học Tự nhiên - Công nghệ} đã truyền đạt cho em nền tảng kiến thức về phân tích hệ thống, cơ sở dữ liệu, phát triển ứng dụng web, kiểm thử phần mềm và triển khai hệ thống. Đó là hành trang quan trọng để em có thể xây dựng đề tài này theo hướng vừa có cơ sở học thuật, vừa có tính thực tiễn.
	
	Em cũng xin cảm ơn gia đình và bạn bè đã luôn động viên, tạo điều kiện thuận lợi trong suốt quá trình học tập và thực hiện đồ án. Sự hỗ trợ về tinh thần và thời gian là nguồn động lực rất lớn để em có thể tập trung hoàn thành công việc đúng định hướng.
	
	Mặc dù đã nỗ lực nghiên cứu, hệ thống và báo cáo vẫn khó tránh khỏi những hạn chế nhất định, đặc biệt trong việc tích hợp cảm biến Internet vạn vật thực tế thay cho dữ liệu mô phỏng, kiểm chứng hiệu quả tự động hóa thiết bị trong môi trường thực địa và tối ưu hóa giao diện trên thiết bị di động. Em rất mong nhận được thêm những góp ý quý báu từ thầy cô và bạn đọc để tiếp tục hoàn thiện đề tài ở các giai đoạn phát triển tiếp theo.
	
	\vspace{30pt}
	
	\begin{flushright}
		\textit{Sơn La, ngày ...... tháng ...... năm 2026}\\[10pt]
		\textit{Tác giả đồ án tốt nghiệp}\\[30pt]
		\textbf{\studentname}
	\end{flushright}
	
	\newpage
	\begin{center}
		\textbf{NHẬN XÉT CỦA GIẢNG VIÊN HƯỚNG DẪN}
	\end{center}
	
	\vspace{1cm}
	
	\begin{enumerate}
		\item Hình thức (bố cục, trình bày, lỗi chính tả, đánh số, hình, bảng, phụ lục)
		
		\vspace{3cm}
		
		\item Nội dung (mục tiêu, phương pháp, kết quả, mức độ hoàn thiện, khả năng mở rộng)
		
		\vspace{3cm}
		
		\item Kết luận
		
		\vspace{3cm}
	\end{enumerate}
	
	\vspace{1cm}
	
	\begin{flushright}
		\textit{Sơn La, Ngày ...... tháng ...... năm 2026}\\[0.5cm]
		\textbf{Giảng viên hướng dẫn}\\[0.3cm]
		\textit{(Ký tên, ghi rõ họ tên)}
	\end{flushright}
	
	\newpage
	\begin{center}
		\textbf{NHẬN XÉT CỦA GIẢNG VIÊN CHẤM}
	\end{center}
	
	\vspace{1cm}
	
	\begin{enumerate}
		\item Hình thức (bố cục, trình bày, lỗi chính tả, đánh số, hình, bảng, phụ lục)
		
		\vspace{3cm}
		
		\item Nội dung (mục tiêu, phương pháp, kết quả, mức độ hoàn thiện, khả năng mở rộng)
		
		\vspace{3cm}
		
		\item Kết luận
		
		\vspace{3cm}
	\end{enumerate}
	
	\vspace{1cm}
	
	\begin{flushright}
		\textit{Sơn La, Ngày ...... tháng ...... năm 2026}\\[0.5cm]
		\textbf{Giảng viên chấm}\\[0.3cm]
		\textit{(Ký tên, ghi rõ họ tên)}
	\end{flushright}
	
	\clearpage
	\tableofcontents
	\clearpage
	
	%	\begin{abstract}
		%		Nông nghiệp ứng dụng công nghệ cao đang trở thành hướng phát triển quan trọng tại các tỉnh miền núi phía Bắc, trong đó có tỉnh Sơn La và khu vực Tây Bắc. Tuy nhiên, nhiều nông hộ và mô hình nhà lưới, nhà kính quy mô vừa và nhỏ vẫn theo dõi môi trường canh tác bằng kinh nghiệm cá nhân hoặc ghi chép thủ công. Cách làm này khiến dữ liệu về nhiệt độ, độ ẩm, độ ẩm đất, ánh sáng, độ pH, độ dẫn điện và nồng độ CO$_2$ bị rời rạc, khó phát hiện sớm bất thường và phản ứng chậm khi điều kiện môi trường vượt ra ngoài ngưỡng phù hợp với cây trồng.
		%		
		%		Đề tài này tập trung phân tích, thiết kế và xây dựng \textbf{Hệ thống giám sát nông nghiệp thông minh} -- một ứng dụng web hỗ trợ số hóa các nghiệp vụ cốt lõi trong quá trình theo dõi và điều khiển môi trường canh tác. Hệ thống phục vụ hai nhóm người dùng chính là quản trị viên và nông dân, đồng thời có một tác vụ nền tự động để thu thập dữ liệu, sinh cảnh báo và kích hoạt thiết bị khi cần can thiệp. Các chức năng chính bao gồm: đăng nhập, tạo hồ sơ cây trồng, tạo khu vực trồng trọt, phân công nông dân vào khu vực, thêm cảm biến, thêm thiết bị thông minh và tự động kích hoạt thiết bị dựa trên dữ liệu cảm biến. Phần máy chủ được xây dựng bằng FastAPI, sử dụng PostgreSQL làm cơ sở dữ liệu; phần giao diện được phát triển với React, TypeScript và Vite; toàn bộ hệ thống được đóng gói bằng Docker Compose để thuận lợi cho triển khai.
		%		
		%		\textbf{Nội dung, phạm vi của đề tài.} Đề tài khảo sát bài toán giám sát môi trường canh tác tại các mô hình nông nghiệp thông minh quy mô vừa và nhỏ, phân tích yêu cầu, thiết kế cơ sở dữ liệu và xây dựng hệ thống web hỗ trợ quản lý hồ sơ cây trồng, khu vực trồng trọt, cảm biến, thiết bị thông minh, cảnh báo và tự động hóa thiết bị. Phạm vi thực hiện tập trung vào phiên bản web phục vụ quản trị viên và nông dân; dữ liệu cảm biến hiện được mô phỏng theo chu kỳ để kiểm chứng luồng nghiệp vụ; cơ chế điều khiển tự động dựa trên ngưỡng môi trường và cấu hình thiết bị liên kết với cảm biến. Phiên bản hiện tại chưa tích hợp trực tiếp với thiết bị Internet vạn vật thực địa, chưa có cơ chế phản hồi khép kín từ cảm biến sau khi bật thiết bị và chưa mở rộng sang dự báo năng suất bằng học máy.
		%	\end{abstract}
	
	\clearpage
	\listoftables
	\clearpage
	\listoffigures
	\clearpage
	
	\begin{center}
		\fontsize{14pt}{17pt}\selectfont
		{DANH MỤC TỪ VIẾT TẮT}
	\end{center}
	
	\vspace{10pt}
	
	\begin{longtable}{|p{3cm}|p{11cm}|}
		\hline
		\textbf{Chữ viết tắt} & \multicolumn{1}{c|}{\textbf{Ý nghĩa}} \\
		\hline
		API & Application Programming Interface, giao diện lập trình ứng dụng giữa phần giao diện và phần máy chủ. \\
		\hline
		BE & Backend, phần xử lý nghiệp vụ, xác thực và truy xuất dữ liệu phía máy chủ. \\
		\hline
		CRUD & Create, Read, Update, Delete, bốn thao tác cơ bản với dữ liệu. \\
		\hline
		DB & Database, cơ sở dữ liệu quan hệ lưu trữ toàn bộ dữ liệu vận hành hệ thống. \\
		\hline
		EC & Electrical Conductivity, độ dẫn điện, dùng để đo nồng độ dinh dưỡng trong dung dịch đất. \\
		\hline
		FE & Frontend, phần giao diện người dùng được xây dựng bằng React và TypeScript. \\
		\hline
		FK & Foreign Key, khóa ngoại trong cơ sở dữ liệu quan hệ. \\
		\hline
		FMIS & Farm Management Information System, hệ thống thông tin quản lý nông trại. \\
		\hline
		HTTPS & Hypertext Transfer Protocol Secure, giao thức truyền tải siêu văn bản an toàn. \\
		\hline
		IoT & Internet of Things, mạng lưới các thiết bị vật lý được kết nối và trao đổi dữ liệu qua Internet. \\
		\hline
		JSON & JavaScript Object Notation, định dạng dữ liệu trao đổi giữa giao diện và máy chủ. \\
		\hline
		JWT & JSON Web Token, chuẩn token dùng cho xác thực người dùng không trạng thái. \\
		\hline
		ORM & Object Relational Mapping, cơ chế ánh xạ đối tượng Python với bảng dữ liệu quan hệ. \\
		\hline
		PA & Precision Agriculture, nông nghiệp chính xác. \\
		\hline
		PAR & Photosynthetically Active Radiation, bức xạ quang hợp hoạt động, đo cường độ ánh sáng hiệu quả cho cây trồng. \\
		\hline
		PK & Primary Key, khóa chính trong cơ sở dữ liệu quan hệ. \\
		\hline
		RBAC & Role-Based Access Control, cơ chế kiểm soát truy cập theo vai trò người dùng. \\
		\hline
		REST & Representational State Transfer, kiến trúc API dựa trên giao thức HTTP. \\
		\hline
		UC & Use Case, ca sử dụng trong mô hình hóa yêu cầu hệ thống. \\
		\hline
		UI & User Interface, giao diện tương tác của người sử dụng. \\
		\hline
		UUID & Universally Unique Identifier, mã định danh duy nhất toàn cục. \\
		\hline
		VPD & Vapor Pressure Deficit, thiếu hụt áp suất hơi nước, chỉ số phản ánh mức độ căng thẳng thoát hơi nước của cây. \\
		\hline
		WSN & Wireless Sensor Network, mạng cảm biến không dây. \\
		\hline
	\end{longtable}
	
	\clearpage
	
	% ============================================================
	% MỞ ĐẦU
	% ============================================================
	\begin{center}
		\section*{MỞ ĐẦU}
	\end{center}
	\addcontentsline{toc}{section}{MỞ ĐẦU}
	
	% -----------------------------------------------------------
	\subsection*{1. Lý do chọn đề tài}
	\addcontentsline{toc}{subsection}{1. Lý do chọn đề tài}
	
	Nông nghiệp ứng dụng công nghệ cao đang trở thành định hướng chiến lược tại nhiều địa phương vùng Tây Bắc Việt Nam. Riêng tỉnh Sơn La, đến năm 2025 đã có hơn 115 ha nhà lưới, nhà kính, hệ thống tưới nhỏ giọt trên hơn 3.200 ha và hơn 5.500 ha đạt chứng nhận VietGAP \cite{boKhoaHoc2025a}. Tuy nhiên, phần lớn nông hộ tại đây vẫn theo dõi môi trường canh tác theo phương thức thủ công: dữ liệu rời rạc, phản ứng chậm khi có sự cố và các lệnh điều khiển thiết bị không được lưu vết, dẫn đến khó khăn trong tối ưu hóa điều kiện sinh trưởng và quản lý hiệu quả vùng trồng.
	
	Trên bình diện quốc tế, các nghiên cứu về nông nghiệp chính xác đã chứng minh rằng việc theo dõi liên tục các thông số môi trường kết hợp cơ chế cảnh báo thông minh có thể tăng năng suất từ 20 đến 30\% đồng thời giảm đáng kể lãng phí tài nguyên \cite{soussi2024}. Các giải pháp thương mại hiện có lại được thiết kế cho nông trại quy mô lớn với chi phí nằm ngoài tầm với của nông hộ vừa và nhỏ; trong khi các prototype sinh viên trong nước thường thiếu bộ quy tắc cảnh báo có căn cứ khoa học và kiến trúc phần mềm đủ chặt chẽ để bảo trì, mở rộng.
	
	Xuất phát từ khoảng trống đó, đề tài \textbf{thiết kế và xây dựng hệ thống giám sát nông nghiệp thông minh } được lựa chọn nhằm xây dựng một nền tảng phần mềm web mã nguồn mở, kiến trúc nhiều lớp, trong đó bộ quy tắc sinh cảnh báo được xây dựng trực tiếp dựa trên ngưỡng tham chiếu khoa học theo từng loại cây trồng \cite{soussi2024}, có cấu hình tự động hóa thiết bị theo cảm biến liên kết và phù hợp với đặc thù vận hành của nông hộ quy mô vừa và nhỏ tại vùng Tây Bắc Việt Nam.
	
	% -----------------------------------------------------------
	\subsection*{2. Mục tiêu nghiên cứu}
	\addcontentsline{toc}{subsection}{2. Mục tiêu nghiên cứu}
	
	Mục tiêu tổng quát của đề tài là xây dựng một nền tảng phần mềm web hỗ trợ người quản lý nông trại và nông hộ theo dõi tình trạng môi trường canh tác theo thời gian gần thực, đồng thời cho phép đưa ra các hành động điều khiển có kiểm soát đối với thiết bị như hệ thống tưới, bơm phân, đèn trồng cây, quạt thông gió và máy sưởi để cải thiện chất lượng của cây trồng.
	
	Cụ thể, đề tài hướng tới các mục tiêu: thiết kế mô hình dữ liệu biểu diễn quan hệ giữa người dùng, khu vực trồng, hồ sơ cây trồng, cảm biến, số đo, thiết bị điều khiển, lệnh tác động và cảnh báo; xây dựng phần máy chủ theo kiến trúc nhiều lớp, tách biệt rõ tầng xử lý và phù hợp với định hướng kiến trúc FMIS dựa trên IoT \cite{koksal2019}; xây dựng giao diện trực quan cho hai vai trò quản trị viên và nông dân; hiện thực bộ quy tắc sinh cảnh báo theo ngưỡng đơn và ngưỡng tổ hợp dựa trên phương pháp luận của Soussi và cộng sự (2024) \cite{soussi2024}; xây dựng cơ chế mô phỏng dữ liệu cảm biến theo chu kỳ để kiểm chứng toàn bộ luồng xử lý; và đóng gói môi trường triển khai bằng Docker để thuận lợi cho vận hành thực tế.
	
	% -----------------------------------------------------------
	\subsection*{3. Đối tượng và phạm vi nghiên cứu}
	\addcontentsline{toc}{subsection}{3. Đối tượng và phạm vi nghiên cứu}
	
	\textbf{Đối tượng nghiên cứu} là bài toán xây dựng hệ thống giám sát nông nghiệp thông minh trên nền tảng web, tập trung vào tổ chức dữ liệu môi trường, chuyển hóa ngưỡng sinh trưởng cây trồng thành logic cảnh báo có căn cứ khoa học và thiết kế giao diện quản trị phù hợp với người dùng nông nghiệp.
	
	\textbf{Phạm vi chức năng:} Xác thực người dùng; quản lý tài khoản nông dân; quản lý hồ sơ cây trồng và khu vực trồng trọt; quản lý cảm biến và chuỗi số liệu; quản lý thiết bị điều khiển, cấu hình tự động hóa và trạng thái vận hành; sinh cảnh báo tự động theo ngưỡng đơn và tổ hợp; kích hoạt thiết bị phù hợp khi có cảnh báo; trang điều khiển tổng hợp.
	
	\textbf{Phạm vi công nghệ:} Phần máy chủ sử dụng FastAPI, SQLAlchemy, Alembic và PostgreSQL; phần giao diện sử dụng React, TypeScript, TanStack Router/Query và Zustand; môi trường triển khai được đóng gói bằng Docker Compose.
	
	\textbf{Giới hạn:} Đề tài chưa tích hợp cổng kết nối IoT thực địa; cơ chế tự động hóa hiện dựa trên bộ quy tắc và thời gian chạy cấu hình sẵn, chưa có phản hồi khép kín từ cảm biến thực sau khi bật thiết bị; hệ thống cũng chưa mở rộng sang học máy dự báo năng suất. Đây là các hướng phát triển tiếp theo được đề xuất sau khi phần lõi phần mềm được ổn định hóa.
	
	% -----------------------------------------------------------
	\subsection*{4. Phương pháp nghiên cứu}
	\addcontentsline{toc}{subsection}{4. Phương pháp nghiên cứu}
	
	Đề tài kết hợp bốn phương pháp nghiên cứu chính.
	
	\textbf{Phương pháp khảo sát và phân tích nghiệp vụ:} Khảo sát quy trình vận hành thực tế tại các mô hình nhà kính và nhà lưới ứng dụng công nghệ cao ở vùng Tây Bắc, từ đó xác định các tác nhân, thực thể dữ liệu và quy trình cần kiểm soát.
	
	\textbf{Phương pháp phân tích tài liệu khoa học:} Nghiên cứu các công trình tổng quan về nông nghiệp chính xác \cite{soussi2024,koksal2019,ojha2015,wolfert2017,goap2018} để định hướng bộ quy tắc cảnh báo và cấu trúc dữ liệu, đảm bảo tính có căn cứ khoa học cho phần logic trọng tâm của hệ thống.
	
	\textbf{Phương pháp thiết kế phần mềm:} Vận dụng kiến trúc nhiều lớp kết hợp UML để đặc tả ca sử dụng, luồng hoạt động và tuần tự tương tác, làm cơ sở triển khai cài đặt.
	
	\textbf{Phương pháp kiểm chứng qua mô phỏng:} Dữ liệu cảm biến được mô phỏng tự động theo chu kỳ với biên độ dao động thực tế để kiểm chứng kiến trúc và bộ quy tắc trước khi tích hợp phần cứng thực địa.
	
	% -----------------------------------------------------------
	\subsection*{5. Cấu trúc báo cáo}
	\addcontentsline{toc}{subsection}{5. Cấu trúc báo cáo}
	
	Cấu trúc bài luận được tổ chức thành ba chương chính:
	
	\textbf{Chương 1 -- Khảo sát và Đặc tả Hệ thống} trình bày tổng quan nghiên cứu về nông nghiệp chính xác, kết quả khảo sát nhu cầu tại vùng Tây Bắc, phân tích các hệ thống liên quan và khoảng trống nghiên cứu, cùng đặc tả đầy đủ tác nhân, yêu cầu chức năng, phi chức năng và ràng buộc nghiệp vụ.
	
	\textbf{Chương 2 -- Phân tích và Thiết kế Hệ thống} trình bày phân tích chi tiết luồng xử lý dữ liệu cảm biến, cơ chế sinh cảnh báo tự động có căn cứ khoa học và cách kích hoạt thiết bị phù hợp khi dữ liệu vượt hoặc thấp hơn ngưỡng. Tiếp theo là mô hình hóa hệ thống qua các ca sử dụng chi tiết kèm biểu đồ hoạt động và tuần tự tương tác, cùng thiết kế cơ sở dữ liệu của hệ thống.
	
	\textbf{Chương 3 -- Xây dựng Chương trình} trình bày kiến trúc chi tiết phần máy chủ và phần giao diện, cơ chế mô phỏng dữ liệu, bộ quy tắc sinh cảnh báo, tự động hóa thiết bị, phương án triển khai Docker Compose và minh họa giao diện hệ thống.
	
	% ============================================================
	% CHƯƠNG 1: KHẢO SÁT VÀ ĐẶC TẢ HỆ THỐNG
	% ============================================================
	
	\chaptersection{KHẢO SÁT VÀ ĐẶC TẢ HỆ THỐNG}
	
	% -----------------------------------------------------------
	\subsection{Tổng quan nghiên cứu}
	
	\subsubsection{Bối cảnh nông nghiệp vùng Tây Bắc và nhu cầu chuyển đổi số}
	
	Vùng Tây Bắc Việt Nam, bao gồm các tỉnh Sơn La, Điện Biên, Lai Châu và Hòa Bình, sở hữu lợi thế địa lý và khí hậu đặc thù với nhiều tiểu vùng khí hậu đan xen, tạo điều kiện thuận lợi cho canh tác đa dạng loài cây từ nhiệt đới đến ôn đới. Riêng tỉnh Sơn La, đặc biệt là cao nguyên Mộc Châu và vùng Mai Sơn, Yên Châu, đã hình thành nhiều vùng sản xuất nông nghiệp hàng hóa tập trung với các sản phẩm đặc trưng như mận hậu, nhãn, xoài, cà phê, chè Shan tuyết và các loại rau ôn đới \cite{cucThuongMai2018}. Đến năm 2025, tỉnh Sơn La đã xây dựng hơn 115 ha nhà lưới, nhà kính, lắp đặt hệ thống tưới nhỏ giọt trên khoảng 3.200 ha cây trồng và hơn 5.500 ha được chứng nhận VietGAP \cite{boKhoaHoc2025a}, từng bước khẳng định định hướng trở thành vùng nông nghiệp ứng dụng công nghệ cao của khu vực.
	
	Tuy nhiên, trong thực tế vận hành, phần lớn các nông hộ tại khu vực này vẫn ghi nhận thông tin canh tác theo phương thức thủ công và quan sát trạng thái khu vực trồng trọt theo định kỳ trước khi đưa ra quyết định tưới nước, bón phân, thông gió hay bổ sung ánh sáng. Cách làm này gặp ba hạn chế lớn có tính hệ thống. Thứ nhất, dữ liệu môi trường rời rạc và khó tổng hợp, dẫn đến khó phát hiện sớm các xu hướng biến động bất thường. Thứ hai, phản ứng với sự cố thường chậm do người vận hành không được cảnh báo kịp thời. Thứ ba, các quyết định điều khiển thiết bị không được lưu vết, gây khó khăn khi đánh giá hiệu quả can thiệp hoặc truy nguyên nguyên nhân sai lệch trong quá trình canh tác.
	
	Theo khảo sát về thực trạng chuyển đổi số nông nghiệp tại các tỉnh miền núi phía Bắc, một trong những rào cản lớn nhất là quy mô sản xuất còn nhỏ lẻ, manh mún, địa hình chia cắt và trình độ tiếp cận công nghệ chưa đồng đều \cite{boKhoaHoc2025b}. Đặc biệt trong bối cảnh Nhà nước đang triển khai các chương trình ứng dụng IoT, tự động hóa và quản trị dữ liệu vào sản xuất nông nghiệp theo định hướng Nghị quyết 57 \cite{boKhoaHoc2025a}, nhu cầu về một nền tảng phần mềm hỗ trợ giám sát và điều khiển môi trường canh tác phù hợp với đặc thù vùng Tây Bắc ngày càng trở nên cấp thiết.
	
	\subsubsection{Tổng quan nghiên cứu về nông nghiệp chính xác và hệ thống giám sát thông minh}
	
	Nông nghiệp chính xác (Precision Agriculture --- PA) đang được thúc đẩy mạnh mẽ bởi sự hội tụ của cảm biến thông minh, phân tích dữ liệu lớn và nền tảng IoT. Nghiên cứu tổng quan của Soussi và cộng sự (2024) chỉ ra rằng mục tiêu cốt lõi của PA là tối ưu hóa thực hành canh tác, nâng cao hiệu quả sử dụng tài nguyên và giảm tác động môi trường thông qua tích hợp cảm biến thông minh với IoT, phân tích dữ liệu lớn và trí tuệ nhân tạo \cite{soussi2024}. Các hệ thống giám sát môi trường vận hành trên nền Mạng Cảm biến Không dây (WSN) cung cấp dịch vụ theo dõi liên tục nhằm duy trì điều kiện tăng trưởng tối ưu và phát hiện sớm các yếu tố có thể gây bệnh hại lan rộng.
	
	Đặc biệt, Soussi và cộng sự (2024) nhấn mạnh vai trò của bảy thông số môi trường đo liên tục --- nhiệt độ không khí, độ ẩm tương đối, độ ẩm đất, cường độ ánh sáng, nồng độ CO$_2$, độ pH và độ dẫn điện EC --- trong việc quản lý sức khỏe cây trồng và hiệu quả sử dụng tài nguyên. Quan trọng hơn, nghiên cứu này đề xuất các ngưỡng tham chiếu cụ thể cho từng thông số theo từng loại cây trồng và điều kiện nhà kính --- đây chính là cơ sở khoa học trực tiếp được đề tài AgriSmart sử dụng để xây dựng bộ quy tắc cảnh báo vượt ngưỡng (được trình bày chi tiết tại Chương 3).
	
	Wolfert và cộng sự (2017) trong nghiên cứu tổng quan về dữ liệu lớn trong nông nghiệp thông minh chỉ ra rằng nông nghiệp hiện đại ngày càng dựa mạnh vào khả năng liên thông dữ liệu và các cơ chế hỗ trợ quyết định theo thời gian gần thực \cite{wolfert2017}. Kết quả phân tích của nhóm tác giả cho thấy giá trị thực sự không nằm ở khối lượng dữ liệu thu thập được, mà ở khả năng chuyển hóa dữ liệu thô thành thông tin hành động cụ thể trong thời gian đủ ngắn để can thiệp kịp thời.
	
	Ojha và cộng sự (2015) chứng minh rằng giá trị thực sự của WSN trong nông nghiệp không chỉ nằm ở khâu đo lường mà ở khả năng kết nối dữ liệu cảm biến với hành động điều khiển và hỗ trợ ra quyết định \cite{ojha2015}. Nghiên cứu này cũng phân tích các ràng buộc kỹ thuật đặc thù của môi trường canh tác --- năng lượng hạn chế, địa hình phức tạp, yêu cầu thời gian thực --- và kết luận rằng thiết kế phần mềm phải đủ linh hoạt để trừu tượng hóa các ràng buộc này khỏi tầng nghiệp vụ.
	
	Ahmed và cộng sự (2018) khẳng định rằng IoT trong nông nghiệp chính xác chỉ phát huy hiệu quả tối đa khi nền tảng phần mềm đủ trưởng thành để chuyển hóa dữ liệu cảm biến thô thành tín hiệu hành động cụ thể \cite{ahmed2018}. Koksal và Tekinerdogan (2019) bổ sung góc nhìn kiến trúc hệ thống: một FMIS (Farm Management Information System) dựa trên IoT cần được tổ chức thành các lớp độc lập --- cảm biến, thu thập, xử lý, trình bày --- để có thể mở rộng và tích hợp phần cứng mới mà không ảnh hưởng đến tầng nghiệp vụ \cite{koksal2019}. Định hướng kiến trúc này được áp dụng trực tiếp trong thiết kế phần máy chủ của AgriSmart.
	
	\subsubsection{Các hệ thống liên quan và khoảng trống nghiên cứu}
	
	Hiện nay trên thị trường đã tồn tại một số nền tảng thương mại phục vụ giám sát và tự động hóa nông nghiệp thông minh. Growlink là hệ thống phần cứng và phần mềm tích hợp của Mỹ, cho phép giám sát và điều khiển từ xa khí hậu, tưới tiêu, bổ sung dinh dưỡng và chiếu sáng trong nhà kính quy mô thương mại lớn \cite{growlink2024}. Farmapp và các nền tảng tương tự trong hệ sinh thái IoT nông nghiệp toàn cầu cung cấp khả năng theo dõi dữ liệu cảm biến thời gian thực kết hợp phân tích học máy để dự báo sản lượng \cite{farmapp2024}. Trong nước, một số đề tài nghiên cứu sinh viên và dự án thí điểm cũng đã xây dựng các bản giám sát nhà kính dựa trên Arduino hoặc Raspberry Pi kết hợp trong tổng quan đơn giản.
	
	Tuy nhiên, khi xem xét trong bối cảnh vùng Tây Bắc Việt Nam, các hệ thống thương mại kể trên bộc lộ hai hạn chế rõ rệt. Thứ nhất, chúng được thiết kế cho quy mô nông trại thương mại lớn tại các nước phát triển, với chi phí triển khai và vận hành nằm ngoài khả năng của hầu hết nông hộ vừa và nhỏ tại miền núi. Thứ hai, bộ quy tắc cảnh báo và logic nghiệp vụ bị đóng gói hoàn toàn bên trong phần cứng độc quyền, không cho phép người dùng tùy chỉnh ngưỡng theo đặc thù từng loại cây trồng địa phương. Đối với các bản mẫu sinh viên trong nước, vấn đề ngược lại: dữ liệu cảm biến được thu thập nhưng chưa có bộ quy tắc sinh cảnh báo dựa trên căn cứ khoa học, chưa có cấu hình liên kết cảm biến với thiết bị phù hợp và kiến trúc phần mềm thường còn đơn lớp, khó bảo trì và mở rộng.
	
	Có thể tóm tắt khoảng trống nghiên cứu qua bảng so sánh sau:
	
	\begin{table}[H]
		\centering
		\caption{So sánh các hệ thống liên quan và định vị AgriSmart}
		\label{tab:related-systems}
		\begin{tabular}{|p{3.5cm}|p{2.5cm}|p{2.5cm}|p{2.5cm}|p{2.5cm}|}
			\hline
			\textbf{Tiêu chí} & \textbf{Growlink / Farmapp} & \textbf{Bản mẫu sinh viên trong nước} & \textbf{AgriSmart} \\
			\hline
			Chi phí triển khai & Cao & Thấp & Thấp \\
			\hline
			Bộ quy tắc cảnh báo có căn cứ khoa học & Có (đóng gói) & Không & Có \\
			\hline
			Lịch sử lệnh điều khiển & Có & Không/sơ sài & Có, đầy đủ \\
			\hline
			Kiến trúc nhiều lớp, dễ mở rộng & Có & Không & Có \\
			\hline
			Phù hợp nông hộ vừa và nhỏ & Không & Có & Có \\
			\hline
		\end{tabular}
	\end{table}
	
	Khoảng trống mà đề tài này hướng đến lấp đầy là: một nền tảng phần mềm web mã nguồn mở, kiến trúc nhiều lớp, chi phí triển khai thấp, trong đó bộ quy tắc sinh cảnh báo được xây dựng trực tiếp trên các ngưỡng tham chiếu khoa học theo từng loại cây trồng, có cấu hình tự động hóa thiết bị theo cảm biến liên kết và phù hợp với quy trình vận hành của nông hộ quy mô vừa và nhỏ tại vùng Tây Bắc Việt Nam.
	
	% -----------------------------------------------------------
	\subsection{Khảo sát thực tế}
	
	\subsubsection{Kết quả khảo sát}
	
	Khảo sát nhu cầu giám sát môi trường canh tác tại các mô hình nhà kính và nhà lưới ứng dụng công nghệ cao ở vùng Tây Bắc, đặc biệt tại Sơn La, cho thấy hoạt động vận hành xoay quanh bốn nhóm nghiệp vụ chính.
	
	\textbf{Giám sát môi trường liên tục:} Người vận hành cần theo dõi liên tục bảy thông số cốt lõi: nhiệt độ không khí, độ ẩm tương đối, độ ẩm đất, cường độ ánh sáng, độ pH đất, độ dẫn điện (EC) và nồng độ CO$_2$ \cite{soussi2024}. Mỗi loại cây trồng có dải ngưỡng tối ưu riêng biệt; khi thông số vượt ngưỡng, cây có nguy cơ bị căng thẳng sinh lý, sâu bệnh hoặc giảm năng suất nghiêm trọng. Việc theo dõi thủ công không đủ tần suất để phát hiện biến động bất thường trong ngày, đặc biệt trong điều kiện khí hậu biến đổi mạnh theo giờ tại vùng miền núi.
	
	\textbf{Quản lý hồ sơ cây trồng và khu vực trồng:} Một nông trại thường có nhiều khu vực trồng khác nhau (nhà kính, nhà lưới, khu ngoài trời), mỗi khu có thể đang canh tác loại cây khác nhau ở các giai đoạn sinh trưởng khác nhau. Người quản lý cần biết khu vực nào đang áp dụng hồ sơ cây trồng nào, ngưỡng tham chiếu nào đang được sử dụng và trạng thái hoạt động của từng khu vực.
	
	\textbf{Quản lý thiết bị thông minh có kiểm soát:} Khi phát hiện thông số bất thường, người vận hành cần kích hoạt hệ thống tưới (điều chỉnh độ ẩm đất), quạt thông gió (điều chỉnh nhiệt độ và CO$_2$), đèn trồng (bổ sung ánh sáng), bơm phân (điều chỉnh EC) hoặc máy sưởi (tăng nhiệt độ khi quá thấp). Thiết bị cần biết mình thuộc khu vực nào, liên kết với cảm biến nào, chạy khi dữ liệu dưới ngưỡng hay vượt ngưỡng và tự tắt sau thời gian đã cấu hình.
	
	\textbf{Quản lý cảnh báo và phân công người vận hành:} Một nông trại có thể có nhiều nông dân phụ trách các khu vực khác nhau. Cần có cơ chế phân quyền để mỗi nông dân chỉ nhận cảnh báo và thao tác trên khu vực của mình, trong khi quản trị viên giám sát toàn hệ thống. Mô hình phân quyền này khác biệt đáng kể so với mô hình owner-based đơn giản thường gặp trong các bản mẫu của sinh viên.
	
	Hiện tại, toàn bộ các hoạt động trên chủ yếu được thực hiện thủ công hoặc dựa vào kinh nghiệm cá nhân, chưa có phần mềm hỗ trợ thu thập dữ liệu liên tục, sinh cảnh báo theo ngưỡng và tự động kích hoạt thiết bị phù hợp một cách có hệ thống.
	
	\subsubsection{Phân tích kết quả khảo sát}
	
	Từ kết quả khảo sát, bài toán giám sát nông nghiệp ẩn chứa sự phức tạp đáng kể khi nhìn từ góc độ luồng dữ liệu và trạng thái nghiệp vụ. Một cảnh báo tưởng chừng đơn giản --- chẳng hạn ``nhiệt độ quá cao'' --- thực tế đòi hỏi hệ thống phải biết: cây trồng nào đang được giám sát, ngưỡng tối ưu của loại cây đó là bao nhiêu, mức độ vượt ngưỡng là bao nhiêu phần trăm (để phân loại thấp, trung bình hoặc cao) và đã có cảnh báo tương tự chưa được xử lý chưa (để tránh lặp cảnh báo liên tục).
	
	Bên cạnh đó, sự phức tạp còn đến từ yêu cầu phát hiện các \textit{điều kiện tổ hợp}: nhiệt độ cao đồng thời với độ ẩm cao là dấu hiệu nguy cơ bệnh nấm, trong khi nhiệt độ cao kết hợp độ ẩm thấp lại gây thiếu hụt áp suất hơi nước. Đây là các tình huống không thể phát hiện bằng cách chỉ theo dõi từng thông số độc lập, dẫn đến yêu cầu thiết kế trọng tâm: hệ thống cần một bộ quy tắc có khả năng đánh giá cả ngưỡng đơn lẫn điều kiện tổ hợp đa thông số.
	
	Phân tích cũng cho thấy sự cần thiết của mô hình phân quyền linh hoạt: một khu vực có thể được phân công cho nhiều nông dân, một nông dân có thể quản lý nhiều khu vực, nhưng quyền truy cập phải được kiểm soát chặt chẽ để đảm bảo mỗi người chỉ thao tác trên phạm vi được phân công. Từ góc độ thiết kế dữ liệu, quan hệ này được biểu diễn bằng một quan hệ phân công trung gian thay vì quan hệ sở hữu trực tiếp.
	
	Tổng hợp lại, khảo sát xác định ba yêu cầu thiết kế cốt lõi mà hệ thống phải đáp ứng: (1) thu thập và lưu trữ dữ liệu cảm biến liên tục theo chuỗi thời gian; (2) đánh giá bộ quy tắc cảnh báo theo ngưỡng đơn và tổ hợp một cách tự động và có căn cứ khoa học; (3) kiểm soát phân quyền truy cập ở mức khu vực, không chỉ ở mức tài khoản.
	
	% -----------------------------------------------------------
	\subsection{Đặc tả hệ thống}
	
	\subsubsection{Tổng quan kiến trúc}
	
	Hệ thống AgriSmart được đặc tả như một ứng dụng web nhiều người dùng, hoạt động theo mô hình máy khách -- máy chủ. Ở phía giao diện, quản trị viên và nông dân thao tác thông qua ứng dụng web được xây dựng bằng React và TypeScript. Ở phía máy chủ, hệ thống tiếp nhận yêu cầu thông qua API RESTful, thực hiện xác thực danh tính, kiểm tra quyền truy cập theo vai trò, áp dụng logic nghiệp vụ và truy xuất dữ liệu từ cơ sở dữ liệu PostgreSQL.
	
	Ngoài các API phục vụ giao diện, hệ thống còn có một tác vụ nền chạy định kỳ mỗi 60 giây: mô phỏng số liệu đo từ cảm biến, lưu vào cơ sở dữ liệu, sau đó đánh giá bộ quy tắc cảnh báo trên dữ liệu vừa thu thập và sinh ra các cảnh báo khi phát hiện vi phạm ngưỡng. Toàn bộ hệ thống được đóng gói bằng Docker với các dịch vụ chính gồm phần máy chủ, phần giao diện, PostgreSQL và Redis.
	
	\subsubsection{Các tác nhân và vai trò}
	
	Hệ thống xác định hai nhóm tác nhân người dùng và một tác nhân hệ thống. \textbf{Quản trị viên} tạo và quản lý tài khoản nông dân, định nghĩa hồ sơ cây trồng với ngưỡng tối ưu cho từng loại, tạo và quản lý khu vực trồng trọt, phân công nông dân vào khu vực và theo dõi tổng quan toàn hệ thống.
	
	\textbf{Nông dân} xem trang tổng quan và số liệu cảm biến của các khu vực được phân công, xem và đánh dấu đã đọc cảnh báo, quản lý thiết bị thông minh trong phạm vi khu vực được phân công, liên kết thiết bị với cảm biến và đặt điều kiện tự động hóa. Mọi yêu cầu truy cập ngoài phạm vi phân công đều bị hệ thống từ chối.
	
	\textbf{Hệ thống nền} là tác nhân tự động chạy mỗi 60 giây để mô phỏng và lưu số liệu cảm biến, đánh giá bộ quy tắc cảnh báo theo ngưỡng của hồ sơ cây trồng đang được áp dụng tại từng khu vực, sinh cảnh báo với cơ chế chống trùng lặp trong vòng 1 giờ và kích hoạt thiết bị phù hợp khi có cấu hình tự động hóa.
	
	\subsubsection{Yêu cầu chức năng}
	
	Hệ thống bao gồm sáu nhóm chức năng chính, được liệt kê dưới đây theo thứ tự từ tầng xác thực đến tầng trình bày.
	
	\begin{itemize}[leftmargin=1.5cm,label=-]
		\item \textbf{Xác thực và phân quyền:} Đăng nhập bằng email/mật khẩu; hệ thống duy trì phiên đăng nhập an toàn trong 7 ngày; kiểm soát truy cập theo hai vai trò quản trị viên và nông dân; đổi mật khẩu; đặt lại mật khẩu qua email.
		\item \textbf{Quản lý hồ sơ cây trồng:} Quản trị viên tạo, sửa, xóa hồ sơ với 7 nhóm ngưỡng môi trường (nhiệt độ không khí, độ ẩm không khí, độ ẩm đất, ánh sáng, pH, EC, CO$_2$); tên hồ sơ phải duy nhất trong hệ thống; ngưỡng bao gồm giá trị tối thiểu, tối ưu và tối đa cho phép.
		\item \textbf{Quản lý khu vực trồng trọt:} Quản trị viên tạo, sửa, xóa khu vực; gán hồ sơ cây trồng cho khu vực (một khu vực -- một hồ sơ tại một thời điểm); phân công nông dân vào khu vực; khi thay đổi hồ sơ, bộ quy tắc tự động áp dụng ngưỡng mới trong chu kỳ tiếp theo.
		\item \textbf{Quản lý cảm biến và số đo:} Thêm, sửa, xóa cảm biến trong khu vực; hệ thống lưu chuỗi số liệu theo thời gian; truy vấn lịch sử đo theo khoảng thời gian và theo cảm biến; dữ liệu được trình bày dạng biểu đồ xu hướng trên trang tổng quan.
		\item \textbf{Quản lý thiết bị thông minh:} Nông dân lọc thiết bị theo khu vực, thêm hoặc cập nhật thiết bị, liên kết thiết bị với cảm biến, chọn điều kiện tự động hóa dưới ngưỡng/vượt ngưỡng/cả hai và đặt thời gian chạy tự động; thiết bị vẫn hỗ trợ bật/tắt thủ công khi cần.
		\item \textbf{Quản lý cảnh báo và tự động hóa:} Hệ thống sinh cảnh báo tự động theo ngưỡng đơn và tổ hợp với ba mức độ nghiêm trọng thấp, trung bình và cao; nông dân xem và đánh dấu đã đọc cảnh báo thuộc khu vực được phân công; quản trị viên xem toàn bộ cảnh báo; cơ chế chống trùng lặp ngăn cùng một loại cảnh báo trong cùng khu vực được sinh lại trong vòng 1 giờ nếu chưa được đọc; thiết bị phù hợp được bật tự động và tự tắt theo thời gian cấu hình.
		\item \textbf{Trang tổng quan:} Thẻ tóm tắt số lượng khu vực, cảm biến, thiết bị và cảnh báo chưa đọc; số liệu cảm biến mới nhất theo khu vực; biểu đồ xu hướng chuỗi thời gian.
	\end{itemize}
	
	\subsubsection{Yêu cầu phi chức năng}
	
	\begin{itemize}[leftmargin=1.5cm,label=-]
		\item \textbf{Bảo mật:} Mật khẩu được băm bằng bcrypt; phiên đăng nhập có thời hạn 7 ngày và được lưu an toàn; mọi yêu cầu đều được kiểm tra xác thực và quyền trước khi xử lý; yêu cầu của nông dân vượt phạm vi khu vực được phân công sẽ bị từ chối.
		\item \textbf{Tính nhất quán dữ liệu:} Các thao tác liên quan đến trạng thái --- cập nhật cấu hình thiết bị, sinh cảnh báo kèm kiểm tra trùng lặp và ghi trạng thái tự động hóa --- được thực hiện trong giao dịch cơ sở dữ liệu để đảm bảo toàn vẹn dữ liệu ngay cả khi có lỗi giữa chừng.
		\item \textbf{Khả năng mở rộng:} Kiến trúc nhiều lớp cho phép bổ sung loại cảm biến mới, loại thiết bị mới hoặc tích hợp cổng kết nối IoT thực mà không ảnh hưởng đến các thành phần hiện có.
		\item \textbf{Khả năng triển khai:} Hệ thống được đóng gói bằng Docker Compose gồm năm dịch vụ, có thể khởi chạy nhất quán trên mọi môi trường với một lệnh duy nhất, không cần cấu hình thủ công từng thành phần.
		\item \textbf{Khả năng kiểm chứng:} Tác vụ nền mô phỏng dữ liệu cảm biến hoạt động độc lập, cho phép kiểm chứng toàn bộ luồng nghiệp vụ từ thu thập số liệu, sinh cảnh báo đến kích hoạt thiết bị tự động mà không cần phần cứng thực địa.
	\end{itemize}
	
	\subsubsection{Ràng buộc nghiệp vụ}
	
	Các ràng buộc nghiệp vụ trọng tâm được đặc tả như sau:
	
	\begin{itemize}[leftmargin=1.5cm,label=-]
		\item Nông dân chỉ được thao tác trên các khu vực được quản trị viên phân công; yêu cầu truy cập khu vực ngoài phạm vi sẽ bị từ chối.
		\item Một khu vực trồng trọt chỉ được gán một hồ sơ cây trồng tại một thời điểm; khi thay đổi hồ sơ, bộ quy tắc cảnh báo tự động áp dụng ngưỡng mới trong chu kỳ tiếp theo.
		\item Cảnh báo có cơ chế chống trùng lặp: cùng một loại cảnh báo cho cùng một khu vực không được sinh lại trong vòng 1 giờ nếu cảnh báo trước chưa được đọc; cảnh báo đã đọc cho phép sinh lại ngay khi điều kiện vi phạm tiếp tục tồn tại.
		\item Thiết bị tự động chỉ được kích hoạt khi thiết bị đang hoạt động, đã bật chế độ tự động hóa, đã liên kết đúng với cảm biến phát sinh cảnh báo và có điều kiện kích hoạt khớp với hướng vi phạm ngưỡng.
		\item Bộ quy tắc cảnh báo đánh giá bảy ngưỡng đơn (tương ứng bảy thông số môi trường) và bốn điều kiện tổ hợp: (i)~nhiệt độ cao + độ ẩm cao $\to$ nguy cơ nấm; (ii)~nhiệt độ cao + độ ẩm thấp $\to$ stress VPD; (iii)~độ ẩm đất cao + pH thấp $\to$ nguy cơ thối rễ; (iv)~CO$_2$ thấp + ánh sáng cao $\to$ lãng phí ánh sáng.
		\item Phân loại mức độ nghiêm trọng dựa trên độ lệch tương đối so với ngưỡng tham chiếu: dưới 10\% là thấp; từ 10 đến 25\% là trung bình; trên 25\% là cao.
	\end{itemize}
	
	\chaptersection{PHÂN TÍCH VÀ THIẾT KẾ HỆ THỐNG}
	
	% -----------------------------------------------------------
	\subsection{Phân tích luồng xử lý dữ liệu, sinh cảnh báo và tự động hóa}
	
	Trước khi bước vào mô hình hóa hệ thống bằng các biểu đồ ca sử dụng và tuần tự, cần phân tích rõ luồng xử lý trung tâm của AgriSmart: cơ chế tự động thu thập dữ liệu cảm biến, đánh giá ngưỡng, sinh cảnh báo và kích hoạt thiết bị phù hợp khi cần can thiệp. Đây là nhân tố kỹ thuật cốt lõi phân biệt hệ thống với các ứng dụng giám sát đơn thuần.
	
	\subsubsection{Tổng quan luồng dữ liệu}
	
	Luồng xử lý được thực hiện bởi một tác vụ nền chạy định kỳ mỗi 60 giây, trải qua năm giai đoạn chính: \textbf{(1) Thu thập dữ liệu} --- mô phỏng hoặc đọc giá trị từ các cảm biến môi trường trong từng khu vực trồng trọt và lưu vào chuỗi thời gian; \textbf{(2) Tra cứu ngưỡng} --- lấy hồ sơ cây trồng đang được áp dụng cho khu vực đó để xác định ngưỡng tối ưu cho từng thông số; \textbf{(3) Đánh giá bộ quy tắc} --- so sánh giá trị thực tế với ngưỡng và đánh giá các điều kiện tổ hợp; \textbf{(4) Sinh cảnh báo} --- tạo cảnh báo kèm mức độ nghiêm trọng và khuyến nghị hành động, với cơ chế chống trùng lặp 1 giờ; \textbf{(5) Kích hoạt tự động hóa} --- tìm thiết bị đã liên kết với cảm biến và bật thiết bị theo cấu hình dưới ngưỡng, vượt ngưỡng hoặc cả hai, sau đó tự tắt theo thời gian chạy đã đặt.
	
	Sơ đồ dưới đây mô tả toàn bộ luồng dữ liệu từ cảm biến đến cảnh báo và thiết bị tự động:
	
	\begin{figure}[H]
		\centering
		\resizebox{\textwidth}{!}{
			\begin{tikzpicture}[
				node distance=0.7cm and 1.4cm,
				box/.style={rectangle, rounded corners=4pt, draw=myblue, fill=lightblue, text width=3.2cm, align=center, minimum height=1cm, font=\normalsize},
				decision/.style={diamond, draw=myblue, fill=lightgreen, text width=2.2cm, align=center, aspect=1.6, font=\normalsize},
				store/.style={rectangle, draw=myblue!60, fill=white, text width=3cm, align=center, minimum height=0.8cm, font=\normalsize, dashed},
				arrow/.style={-{Stealth[length=5pt]}, myblue, thick},
				label/.style={font=\normalsize, color=linegray}
				]
				% Nodes
				\node[box] (sensor) {Cảm biến\\(7 loại)};
				\node[box, right=of sensor] (collect) {Thu thập \&\\Lưu số liệu};
				
				\node[box, right=of collect] (profile) {Tra cứu\\Hồ sơ cây trồng};
				
				\node[box, right=of profile] (engine) {Bộ quy tắc\\Đánh giá ngưỡng};
				\node[decision, right=of engine] (vipham) {Vi phạm\\ngưỡng?};
				\node[box, above right=0.4cm and 0.8cm of vipham] (dedup) {Kiểm tra\\trùng lặp};
				\node[store, above=0.4cm of dedup] (alerts_db) {Cảnh báo};
				\node[box, below right=0.4cm and 0.8cm of vipham] (skip) {Bỏ qua\\chu kỳ này};
				\node[box, right=1.2cm of dedup] (severity) {Phân mức\\nghiêm trọng};
				\node[box, right=of severity] (alert_out) {Sinh Cảnh Báo\\+ Khuyến nghị};
				\node[box, below=0.8cm of alert_out] (auto_device) {Kích hoạt\\thiết bị phù hợp};
				\node[store, right=0.6cm of auto_device] (devices_db) {Thiết bị};
				
				% Arrows
				\draw[arrow] (sensor) -- (collect);
				\draw[arrow] (collect) -- (profile);
				
				\draw[arrow] (profile) -- (engine);
				\draw[arrow] (engine) -- (vipham);
				\draw[arrow] (vipham) -- node[label, above left]{Có} (dedup);
				\draw[arrow] (vipham) -- node[label, below left]{Không} (skip);
				\draw[arrow] (dedup) -- node[label, above]{Chưa có} (severity);
				\draw[arrow] (dedup.north) -- (alerts_db.south);
				\draw[arrow] (severity) -- (alert_out);
				\draw[arrow] (alert_out) -- ++(0,-1.5) -| (alerts_db);
				\draw[arrow] (alert_out) -- (auto_device);
				\draw[arrow] (auto_device) -- (devices_db);
			\end{tikzpicture}
		}
		\caption{Luồng xử lý dữ liệu, sinh cảnh báo và kích hoạt tự động hóa}
		\label{fig:data-flow-alert}
	\end{figure}
	
	\subsubsection{Chi tiết từng giai đoạn}
	
	\textbf{Giai đoạn 1 --- Thu thập dữ liệu cảm biến.} Mỗi 60 giây, tác vụ nền duyệt qua toàn bộ khu vực trồng đang hoạt động. Với mỗi cảm biến trong khu vực, hệ thống mô phỏng giá trị đo dựa trên khoảng dao động thực tế theo từng loại cảm biến (nhiệt độ, độ ẩm không khí, độ ẩm đất, ánh sáng, pH, EC, CO$_2$) và lưu số liệu kèm thời điểm đo vào cơ sở dữ liệu. Khi tích hợp phần cứng thực, giai đoạn này được thay bằng lệnh đọc từ cổng kết nối IoT mà không ảnh hưởng đến các bước tiếp theo.
	
	\textbf{Giai đoạn 2 --- Tra cứu ngưỡng theo hồ sơ cây trồng.} Sau khi lưu số liệu, bộ quy tắc tra cứu hồ sơ cây trồng đang được gán cho khu vực đó. Hồ sơ lưu trữ các cặp ngưỡng tối thiểu và tối đa cho từng thông số môi trường, được thiết lập theo căn cứ khoa học từ nghiên cứu Soussi và cộng sự (2024) \cite{soussi2024}. Nếu khu vực chưa được gán hồ sơ, bước đánh giá bị bỏ qua.
	
	\textbf{Giai đoạn 3 --- Đánh giá bộ quy tắc.} Bộ quy tắc thực hiện hai lớp đánh giá đồng thời. Với \textbf{ngưỡng đơn}, hệ thống so sánh từng thông số với ngưỡng tối thiểu và ngưỡng tối đa của hồ sơ cây trồng; nếu vượt ngưỡng, hệ thống tính độ lệch tương đối so với ngưỡng tham chiếu và phân mức cảnh báo thành thấp, trung bình hoặc cao.
	
	Với \textbf{điều kiện tổ hợp}, hệ thống phát hiện các kết hợp thông số nguy hiểm mà đánh giá đơn lẻ không thể nhận ra, gồm: nhiệt độ cao kết hợp độ ẩm cao gây nguy cơ nấm bệnh; nhiệt độ cao kết hợp độ ẩm thấp gây căng thẳng sinh lý do thiếu hụt áp suất hơi nước; độ ẩm đất cao kết hợp pH thấp gây nguy cơ thối rễ; và CO$_2$ thấp kết hợp ánh sáng cao gây lãng phí quang hợp.
	
	Ý nghĩa khoa học của từng điều kiện tổ hợp được giải thích cụ thể: khi nhiệt độ vượt 30\textdegree C đồng thời độ ẩm vượt 85\%, môi trường tạo điều kiện lý tưởng cho nấm phát triển và lan rộng nhanh chóng; khi nhiệt độ cao kết hợp độ ẩm thấp, chỉ số VPD tăng cao khiến cây mất nước qua lá nhanh hơn khả năng hút nước từ rễ, dẫn đến héo và giảm năng suất.
	
	\textbf{Giai đoạn 4 --- Kiểm tra trùng lặp và sinh cảnh báo.} Trước khi ghi cảnh báo mới, hệ thống kiểm tra: nếu đã tồn tại một cảnh báo \textit{cùng loại, cùng khu vực, cùng thông số} được tạo trong vòng 1 giờ gần nhất và \textit{chưa được đọc}, cảnh báo mới sẽ bị bỏ qua để tránh lặp thông báo liên tục. Ngược lại, nếu cảnh báo cũ đã được đọc, hệ thống cho phép sinh lại ngay khi điều kiện vi phạm vẫn tiếp diễn --- phản ánh đúng thực tế người vận hành đã xử lý nhưng vấn đề chưa được khắc phục. Mỗi cảnh báo được lưu kèm giá trị thực tế, ngưỡng vi phạm, mức độ nghiêm trọng và khuyến nghị hành động cụ thể theo từng mức.
	
	\textbf{Giai đoạn 5 --- Kích hoạt tự động hóa thiết bị.} Với cảnh báo ngưỡng đơn, hệ thống xác định hướng vi phạm: giá trị thấp hơn ngưỡng tối thiểu hoặc cao hơn ngưỡng tối đa. Sau đó, hệ thống tìm thiết bị đang hoạt động, đã bật tự động hóa, có liên kết với đúng cảm biến và có điều kiện kích hoạt phù hợp. Ví dụ: nhiệt độ thấp thì thiết bị có cấu hình ``dưới ngưỡng'' như máy sưởi được bật; nhiệt độ cao thì thiết bị có cấu hình ``vượt ngưỡng'' như quạt thông gió được bật. Thiết bị chỉ chạy trong khoảng thời gian đã cấu hình để tránh bật quá lâu gây lãng phí hoặc làm môi trường bị điều chỉnh quá mức.
	
	\subsubsection{Phân tích tác động đến cây trồng}
	
	Giá trị cốt lõi của luồng xử lý trên không chỉ nằm ở việc phát hiện sự cố mà còn ở việc \textit{hướng dẫn hành động} cụ thể để cải thiện điều kiện sinh trưởng. Bảng dưới đây minh họa cách hệ thống chuyển hóa một vi phạm ngưỡng thành khuyến nghị hành động phân theo mức độ:
	
	\begin{table}[H]
		\centering
		\caption{Ví dụ phân tích cảnh báo và khuyến nghị hành động theo mức độ nghiêm trọng}
		\label{tab:alert-action}
		\begin{tabular}{|p{2.5cm}|p{2cm}|p{4cm}|p{5.5cm}|}
			\hline
			\textbf{Thông số} & \textbf{Mức độ} & \textbf{Tình huống} & \textbf{Khuyến nghị hành động} \\
			\hline
			Nhiệt độ cao & Thấp & Lệch ngưỡng $< 10\%$ & Theo dõi thêm; kiểm tra thông gió. \\
			\hline
			Nhiệt độ cao & Trung bình & Lệch $10-25\%$ & Bật quạt thông gió; giảm thời lượng chiếu sáng nhân tạo. \\
			\hline
			Nhiệt độ cao & Cao & Lệch $> 25\%$ & Bật quạt ngay; mở cửa thông gió; giảm tải nhiệt. \\
			\hline
			Độ ẩm đất thấp & Trung bình & Đất khô, ngưỡng lệch 15\% & Kích hoạt hệ thống tưới nhỏ giọt. \\
			\hline
			CO$_2$ thấp + ánh sáng cao & Tổ hợp & Quang hợp kém hiệu quả & Giảm cường độ đèn; kiểm tra thông gió để tăng CO$_2$. \\
			\hline
			Nhiệt độ cao + Độ ẩm cao & Tổ hợp & Nguy cơ nấm bệnh & Bật quạt thông gió; giảm tưới; theo dõi dấu hiệu nấm. \\
			\hline
		\end{tabular}
	\end{table}
	
	Như vậy, thay vì chỉ thông báo ``nhiệt độ quá cao'', hệ thống cung cấp đủ ngữ cảnh để người vận hành ra quyết định can thiệp phù hợp ngay lập tức --- góp phần giảm thiểu thiệt hại cho cây trồng và tối ưu hóa điều kiện sinh trưởng theo thời gian.
	
	% -----------------------------------------------------------
	\subsection{Mô hình hóa hệ thống}
	
	\begin{figure}[H]
		\centering
		\safeincludegraphics[width=0.7\linewidth]{uc_tong_quat}
		\caption{Sơ đồ ca sử dụng tổng quát}
		\label{fig:uctongquat}
	\end{figure}
	Hình \ref{fig:uctongquat} mô tả tổng quát các ca sử dụng chính của hệ thống AgriSmart với ba tác nhân tham gia vào quá trình vận hành: \textbf{quản trị viên}, \textbf{nông dân} và \textbf{hệ thống nền}. Quản trị viên thiết lập hồ sơ cây trồng, tạo khu vực trồng trọt và phân công nông dân vào từng khu vực. Nông dân là nhóm người dùng vận hành trực tiếp, có thể thêm cảm biến và thiết bị thông minh trong các khu vực được phân công. Hệ thống nền chạy định kỳ để đọc dữ liệu cảm biến, kiểm tra ngưỡng và tự động kích hoạt thiết bị phù hợp khi có vi phạm.
	
	Hệ thống được tổ chức thành bảy ca sử dụng chính (UC01--UC07): đăng nhập (UC01), tạo hồ sơ cây trồng (UC02), tạo khu vực trồng trọt (UC03), phân công nông dân vào khu vực (UC04), thêm cảm biến (UC05), thêm thiết bị thông minh (UC06) và tự động kích hoạt thiết bị thông minh (UC07). Cách chia này tập trung vào các thao tác cốt lõi để thiết lập vùng trồng, thu thập dữ liệu và điều khiển tự động.
	
	\subsubsection{Ca sử dụng đăng nhập}
	
	\paragraph{Biểu đồ ca sử dụng}
	\noindent
	\begin{figure}[H]
		\centering
		\safeincludegraphics[width=0.85\linewidth]{images/uc_uc01}
		\caption{Biểu đồ ca sử dụng đăng nhập}
		\label{fig:ucuc01}
	\end{figure}
	Hình \ref{fig:ucuc01} mô tả ca sử dụng đăng nhập với hai nhóm người dùng là quản trị viên và nông dân. Người dùng nhập thông tin tài khoản, hệ thống kiểm tra tính hợp lệ và cho phép truy cập vào các chức năng phù hợp với vai trò.
	
	\begin{table}[H]
		\centering
		\caption{Đặc tả ca sử dụng đăng nhập}
		\label{tab:usecase-dangnhap}
		\begin{tabular}{|p{4cm}|p{10cm}|}
			\hline
			\textbf{Tên ca sử dụng} & \textbf{Đăng nhập hệ thống} \\
			\hline
			Tác nhân chính & Quản trị viên và nông dân \\
			\hline
			Dòng sự kiện chính &
			\begin{minipage}[t]{10cm}
				\vspace{2pt}
				\begin{enumerate}
					\item Người dùng nhập email và mật khẩu trên màn hình đăng nhập.
					\item Hệ thống xác thực thông tin đăng nhập và kiểm tra trạng thái tài khoản.
					\item Nếu hợp lệ, hệ thống tạo phiên đăng nhập an toàn cho người dùng.
					\item Hệ thống điều hướng đến trang tổng quan tương ứng với vai trò người dùng.
				\end{enumerate}
				\vspace{2pt}
			\end{minipage} \\
			\hline
			Dòng sự kiện phụ & Nếu thông tin sai hoặc tài khoản không hoạt động, hệ thống hiển thị thông báo lỗi phù hợp. Tính năng quên mật khẩu cho phép người dùng đặt lại qua email xác thực có thời hạn. \\
			\hline
			Điều kiện tiên quyết & Người dùng đã có tài khoản được quản trị viên tạo trong hệ thống. \\
			\hline
			Điều kiện sau & Người dùng đăng nhập thành công và có quyền truy cập các chức năng tương ứng với vai trò. \\
			\hline
		\end{tabular}
	\end{table}
	
	\paragraph{Biểu đồ hoạt động}
	\noindent
	\begin{figure}[H]
		\centering
		\safeincludegraphics[width=0.45\linewidth]{images/activity_uc01}
		\caption{Biểu đồ hoạt động đăng nhập}
		\label{fig:activityuc01}
	\end{figure}
	Hình \ref{fig:activityuc01} mô tả các bước người dùng nhập thông tin đăng nhập, hệ thống kiểm tra tài khoản, sau đó cho phép truy cập nếu hợp lệ hoặc hiển thị lỗi nếu không hợp lệ.
	
	\paragraph{Biểu đồ tuần tự}
	\noindent
	\begin{figure}[H]
		\centering
		\safeincludegraphics[width=0.9\linewidth]{images/sequence_uc01}
		\caption{Biểu đồ tuần tự đăng nhập}
		\label{fig:sequenceuc01}
	\end{figure}
	Hình \ref{fig:sequenceuc01} mô tả quá trình giao diện gửi thông tin đăng nhập đến máy chủ. Máy chủ kiểm tra tài khoản trong cơ sở dữ liệu, tạo phiên đăng nhập khi hợp lệ và trả kết quả để giao diện điều hướng người dùng vào hệ thống.
	
	% -------------------------------------------------------
	\subsubsection{Ca sử dụng tạo hồ sơ cây trồng}
	
	\paragraph{Biểu đồ ca sử dụng}
	\noindent
	\begin{figure}[H]
		\centering
		\safeincludegraphics[width=0.85\linewidth]{images/uc_uc02}
		\caption{Biểu đồ ca sử dụng tạo hồ sơ cây trồng}
		\label{fig:ucuc02}
	\end{figure}
	Hình \ref{fig:ucuc02} mô tả ca sử dụng tạo hồ sơ cây trồng. Quản trị viên nhập thông tin cây trồng và các ngưỡng môi trường phù hợp để hệ thống dùng làm căn cứ cảnh báo và tự động hóa.
	
	\begin{table}[H]
		\centering
		\caption{Đặc tả ca sử dụng tạo hồ sơ cây trồng}
		\label{tab:usecase-taohosocaytrong}
		\begin{tabular}{|p{4cm}|p{10cm}|}
			\hline
			\textbf{Tên ca sử dụng} & \textbf{Tạo hồ sơ cây trồng} \\
			\hline
			Tác nhân chính & Quản trị viên \\
			\hline
			Dòng sự kiện chính &
			\begin{minipage}[t]{10cm}
				\vspace{2pt}
				\begin{enumerate}
					\item Quản trị viên mở màn hình hồ sơ cây trồng.
					\item Quản trị viên nhập tên cây trồng, mô tả và các ngưỡng môi trường cần theo dõi.
					\item Hệ thống kiểm tra dữ liệu hợp lệ và tên hồ sơ chưa bị trùng.
					\item Hệ thống lưu hồ sơ cây trồng vào cơ sở dữ liệu.
					\item Các khu vực dùng hồ sơ này sẽ được bộ quy tắc cảnh báo đánh giá theo ngưỡng vừa lưu.
				\end{enumerate}
				\vspace{2pt}
			\end{minipage} \\
			\hline
			Dòng sự kiện phụ & Nếu tên hồ sơ đã tồn tại hoặc ngưỡng nhập không hợp lệ, hệ thống hiển thị thông báo lỗi và không lưu dữ liệu. \\
			\hline
			Điều kiện tiên quyết & Người dùng đã đăng nhập với vai trò quản trị viên. \\
			\hline
			Điều kiện sau & Hồ sơ cây trồng được tạo và có thể gán cho khu vực trồng trọt. \\
			\hline
		\end{tabular}
	\end{table}
	
	\paragraph{Biểu đồ hoạt động}
	\noindent
	\begin{figure}[H]
		\centering
		\safeincludegraphics[width=0.45\linewidth]{images/activity_uc02}
		\caption{Biểu đồ hoạt động tạo hồ sơ cây trồng}
		\label{fig:activityuc02}
	\end{figure}
	Hình \ref{fig:activityuc02} mô tả các bước quản trị viên nhập thông tin hồ sơ cây trồng. Hệ thống kiểm tra dữ liệu và tên hồ sơ, nếu hợp lệ thì lưu hồ sơ để dùng cho các lần đánh giá ngưỡng tiếp theo.
	
	\paragraph{Biểu đồ tuần tự}
	\noindent
	\begin{figure}[H]
		\centering
		\safeincludegraphics[width=0.9\linewidth]{images/sequence_uc02}
		\caption{Biểu đồ tuần tự tạo hồ sơ cây trồng}
		\label{fig:sequenceuc02}
	\end{figure}
	Hình \ref{fig:sequenceuc02} mô tả quá trình giao diện gửi thông tin hồ sơ cây trồng đến máy chủ. Máy chủ kiểm tra dữ liệu, lưu hồ sơ vào cơ sở dữ liệu và trả kết quả tạo mới cho giao diện.
	
	% -------------------------------------------------------
	\subsubsection{Ca sử dụng tạo khu vực trồng trọt}
	
	\paragraph{Biểu đồ ca sử dụng}
	\noindent
	\begin{figure}[H]
		\centering
		\safeincludegraphics[width=0.85\linewidth]{images/uc_uc03}
		\caption{Biểu đồ ca sử dụng tạo khu vực trồng trọt}
		\label{fig:ucuc03}
	\end{figure}
	Hình \ref{fig:ucuc03} mô tả ca sử dụng tạo khu vực trồng trọt. Quản trị viên khai báo thông tin khu vực, chọn hồ sơ cây trồng áp dụng và lưu khu vực để phục vụ phân công, lắp cảm biến và theo dõi cảnh báo.
	
	\begin{table}[H]
		\centering
		\caption{Đặc tả ca sử dụng tạo khu vực trồng trọt}
		\label{tab:usecase-taokhuvuc}
		\begin{tabular}{|p{4cm}|p{10cm}|}
			\hline
			\textbf{Tên ca sử dụng} & \textbf{Tạo khu vực trồng trọt} \\
			\hline
			Tác nhân chính & Quản trị viên \\
			\hline
			Dòng sự kiện chính &
			\begin{minipage}[t]{10cm}
				\vspace{2pt}
				\begin{enumerate}
					\item Quản trị viên mở màn hình quản lý khu vực trồng trọt.
					\item Quản trị viên nhập tên khu vực, vị trí, diện tích và ngày trồng nếu có.
					\item Quản trị viên chọn hồ sơ cây trồng áp dụng cho khu vực.
					\item Hệ thống kiểm tra dữ liệu hợp lệ.
					\item Hệ thống tạo khu vực mới và đặt trạng thái hoạt động.
					\item Khu vực mới xuất hiện trong danh sách để phân công nông dân và lắp cảm biến.
				\end{enumerate}
				\vspace{2pt}
			\end{minipage} \\
			\hline
			Dòng sự kiện phụ & Nếu hồ sơ cây trồng không tồn tại hoặc dữ liệu nhập thiếu thông tin cần thiết, hệ thống thông báo lỗi. \\
			\hline
			Điều kiện tiên quyết & Quản trị viên đã đăng nhập; hồ sơ cây trồng đã được tạo nếu muốn gán ngay cho khu vực. \\
			\hline
			Điều kiện sau & Khu vực trồng trọt được lưu và có thể dùng trong phân công, cảm biến, cảnh báo và tự động hóa. \\
			\hline
		\end{tabular}
	\end{table}
	
	\paragraph{Biểu đồ hoạt động}
	\noindent
	\begin{figure}[H]
		\centering
		\safeincludegraphics[width=0.45\linewidth]{images/activity_uc03}
		\caption{Biểu đồ hoạt động tạo khu vực trồng trọt}
		\label{fig:activityuc03}
	\end{figure}
	Hình \ref{fig:activityuc03} mô tả các bước quản trị viên nhập thông tin khu vực và chọn hồ sơ cây trồng. Hệ thống kiểm tra dữ liệu, nếu hợp lệ thì tạo khu vực mới trong cơ sở dữ liệu.
	
	\paragraph{Biểu đồ tuần tự}
	\noindent
	\begin{figure}[H]
		\centering
		\safeincludegraphics[width=0.9\linewidth]{images/sequence_uc03}
		\caption{Biểu đồ tuần tự tạo khu vực trồng trọt}
		\label{fig:sequenceuc03}
	\end{figure}
	Hình \ref{fig:sequenceuc03} mô tả quá trình giao diện gửi thông tin khu vực trồng trọt đến máy chủ. Máy chủ kiểm tra quyền của quản trị viên, lưu thông tin khu vực vào cơ sở dữ liệu và trả kết quả tạo mới cho giao diện.
	
	% -------------------------------------------------------
	\subsubsection{Ca sử dụng phân công nông dân vào khu vực}
	
	\paragraph{Biểu đồ ca sử dụng}
	\noindent
	\begin{figure}[H]
		\centering
		\safeincludegraphics[width=0.85\linewidth]{images/uc_uc04}
		\caption{Biểu đồ ca sử dụng phân công nông dân vào khu vực}
		\label{fig:ucuc04}
	\end{figure}
	Hình \ref{fig:ucuc04} mô tả ca sử dụng phân công nông dân vào khu vực. Quản trị viên chọn khu vực và nông dân phụ trách, hệ thống kiểm tra thông tin rồi cấp quyền thao tác cho nông dân trong khu vực đó.
	
	\begin{table}[H]
		\centering
		\caption{Đặc tả ca sử dụng phân công nông dân vào khu vực}
		\label{tab:usecase-phancongnongdan}
		\begin{tabular}{|p{4cm}|p{10cm}|}
			\hline
			\textbf{Tên ca sử dụng} & \textbf{Phân công nông dân vào khu vực} \\
			\hline
			Tác nhân chính & Quản trị viên \\
			\hline
			Dòng sự kiện chính &
			\begin{minipage}[t]{10cm}
				\vspace{2pt}
				\begin{enumerate}
					\item Quản trị viên chọn khu vực trồng trọt cần phân công.
					\item Hệ thống hiển thị danh sách nông dân hiện có.
					\item Quản trị viên chọn nông dân cần gán vào khu vực.
					\item Hệ thống kiểm tra khu vực và nông dân tồn tại, đồng thời kiểm tra phân công chưa bị trùng.
					\item Hệ thống lưu bản ghi phân công.
					\item Nông dân được cấp quyền xem dữ liệu, thêm cảm biến và thêm thiết bị trong khu vực đó.
				\end{enumerate}
				\vspace{2pt}
			\end{minipage} \\
			\hline
			Dòng sự kiện phụ & Nếu nông dân đã được phân công vào khu vực, hệ thống thông báo trùng phân công. Quản trị viên có thể gỡ phân công khi nông dân không còn phụ trách khu vực đó. \\
			\hline
			Điều kiện tiên quyết & Quản trị viên đã đăng nhập; tài khoản nông dân và khu vực trồng trọt đã tồn tại. \\
			\hline
			Điều kiện sau & Cơ sở dữ liệu được cập nhật và quyền truy cập khu vực của nông dân có hiệu lực ngay. \\
			\hline
		\end{tabular}
	\end{table}
	
	\paragraph{Biểu đồ hoạt động}
	\noindent
	\begin{figure}[H]
		\centering
		\safeincludegraphics[width=0.65\linewidth]{images/activity_uc04}
		\caption{Biểu đồ hoạt động phân công nông dân vào khu vực}
		\label{fig:activityuc04}
	\end{figure}
	Hình \ref{fig:activityuc04} mô tả các bước quản trị viên chọn khu vực và nông dân. Hệ thống kiểm tra dữ liệu, nếu chưa có phân công thì lưu thông tin và cấp quyền truy cập khu vực cho nông dân.
	
	\paragraph{Biểu đồ tuần tự}
	\noindent
	\begin{figure}[H]
		\centering
		\safeincludegraphics[width=0.9\linewidth]{images/sequence_uc04}
		\caption{Biểu đồ tuần tự phân công nông dân vào khu vực}
		\label{fig:sequenceuc04}
	\end{figure}
	Hình \ref{fig:sequenceuc04} mô tả quá trình giao diện gửi yêu cầu phân công đến máy chủ. Máy chủ kiểm tra quyền quản trị viên, kiểm tra thông tin khu vực và nông dân, sau đó lưu phân công vào cơ sở dữ liệu.
	
	% -------------------------------------------------------
	\subsubsection{Ca sử dụng thêm cảm biến}
	
	\paragraph{Biểu đồ ca sử dụng}
	\noindent
	\begin{figure}[H]
		\centering
		\safeincludegraphics[width=0.85\linewidth]{images/uc_uc05}
		\caption{Biểu đồ ca sử dụng thêm cảm biến}
		\label{fig:ucuc05}
	\end{figure}
	Hình \ref{fig:ucuc05} mô tả ca sử dụng thêm cảm biến. Nông dân chọn khu vực được phân công, nhập thông tin cảm biến và hệ thống kiểm tra quyền trước khi lưu cảm biến vào khu vực đó.
	
	\begin{table}[H]
		\centering
		\caption{Đặc tả ca sử dụng thêm cảm biến}
		\label{tab:usecase-themcambien}
		\begin{tabular}{|p{4cm}|p{10cm}|}
			\hline
			\textbf{Tên ca sử dụng} & \textbf{Thêm cảm biến} \\
			\hline
			Tác nhân chính & Nông dân \\
			\hline
			Dòng sự kiện chính &
			\begin{minipage}[t]{10cm}
				\vspace{2pt}
				\begin{enumerate}
					\item Người dùng mở màn hình quản lý cảm biến.
					\item Người dùng chọn khu vực trồng trọt cần lắp cảm biến.
					\item Người dùng nhập tên cảm biến, loại cảm biến, vị trí lắp và chu kỳ cập nhật.
					\item Hệ thống kiểm tra người dùng có quyền trên khu vực.
					\item Hệ thống lưu cảm biến vào cơ sở dữ liệu.
					\item Cảm biến mới được dùng trong luồng thu thập dữ liệu và đánh giá ngưỡng.
				\end{enumerate}
				\vspace{2pt}
			\end{minipage} \\
			\hline
			Dòng sự kiện phụ & Nếu người dùng không có quyền trên khu vực, hệ thống từ chối thao tác. Nếu không nhập đơn vị đo, hệ thống tự chọn đơn vị theo loại cảm biến. \\
			\hline
			Điều kiện tiên quyết & Khu vực trồng trọt đã tồn tại; nông dân chỉ thao tác được trên khu vực đã được phân công. \\
			\hline
			Điều kiện sau & Cảm biến được lưu trong cơ sở dữ liệu và sẵn sàng sinh dữ liệu đo. \\
			\hline
		\end{tabular}
	\end{table}
	
	\paragraph{Biểu đồ hoạt động}
	\noindent
	\begin{figure}[H]
		\centering
		\safeincludegraphics[width=0.45\linewidth]{images/activity_uc05}
		\caption{Biểu đồ hoạt động thêm cảm biến}
		\label{fig:activityuc05}
	\end{figure}
	Hình \ref{fig:activityuc05} mô tả các bước người dùng chọn khu vực và nhập thông tin cảm biến. Hệ thống kiểm tra quyền khu vực và lưu cảm biến nếu dữ liệu hợp lệ.
	
	\paragraph{Biểu đồ tuần tự}
	\noindent
	\begin{figure}[H]
		\centering
		\safeincludegraphics[width=0.9\linewidth]{images/sequence_uc05}
		\caption{Biểu đồ tuần tự thêm cảm biến}
		\label{fig:sequenceuc05}
	\end{figure}
	Hình \ref{fig:sequenceuc05} mô tả quá trình giao diện gửi thông tin cảm biến đến máy chủ. Máy chủ kiểm tra quyền thao tác trên khu vực, lưu cảm biến vào cơ sở dữ liệu và trả kết quả tạo mới cho giao diện.
	
	% -------------------------------------------------------
	\subsubsection{Ca sử dụng thêm thiết bị thông minh}
	
	\paragraph{Biểu đồ ca sử dụng}
	\noindent
	\begin{figure}[H]
		\centering
		\safeincludegraphics[width=0.85\linewidth]{images/uc_uc06}
		\caption{Biểu đồ ca sử dụng thêm thiết bị thông minh}
		\label{fig:ucuc06}
	\end{figure}
	Hình \ref{fig:ucuc06} mô tả ca sử dụng thêm thiết bị thông minh. Nông dân khai báo thiết bị, liên kết thiết bị với cảm biến và thiết lập điều kiện bật tự động khi dữ liệu cảm biến thấp hơn hoặc cao hơn ngưỡng.
	
	\begin{table}[H]
		\centering
		\caption{Đặc tả ca sử dụng thêm thiết bị thông minh}
		\label{tab:usecase-themthietbi}
		\begin{tabular}{|p{4cm}|p{10cm}|}
			\hline
			\textbf{Tên ca sử dụng} & \textbf{Thêm thiết bị thông minh} \\
			\hline
			Tác nhân chính & Nông dân \\
			\hline
			Dòng sự kiện chính &
			\begin{minipage}[t]{10cm}
				\vspace{2pt}
				\begin{enumerate}
					\item Nông dân mở màn hình thiết bị thông minh.
					\item Nông dân nhập tên, vị trí, loại thiết bị và thông tin kết nối.
					\item Nông dân chọn cảm biến liên kết để biết thiết bị phục vụ khu vực nào.
					\item Nông dân bật hoặc tắt chế độ tự động hóa cho thiết bị.
					\item Nông dân chọn điều kiện kích hoạt: dưới ngưỡng, vượt ngưỡng hoặc cả hai.
					\item Nông dân nhập thời gian chạy trước khi thiết bị tự tắt.
					\item Hệ thống kiểm tra quyền với cảm biến liên kết và lưu thiết bị.
				\end{enumerate}
				\vspace{2pt}
			\end{minipage} \\
			\hline
			Dòng sự kiện phụ & Nếu cảm biến liên kết không thuộc khu vực nông dân được phân công, hệ thống từ chối thao tác. Nếu tắt tự động hóa, thiết bị chỉ dùng cho điều khiển thủ công. \\
			\hline
			Điều kiện tiên quyết & Nông dân đã được phân công vào khu vực có cảm biến cần liên kết. \\
			\hline
			Điều kiện sau & Thiết bị được lưu trong bảng cơ sở dữ liệu và có thể được dùng bởi luồng tự động hóa. \\
			\hline
		\end{tabular}
	\end{table}
	
	\paragraph{Biểu đồ hoạt động}
	\noindent
	\begin{figure}[H]
		\centering
		\safeincludegraphics[width=0.35\linewidth]{images/activity_uc06}
		\caption{Biểu đồ hoạt động thêm thiết bị thông minh}
		\label{fig:activityuc06}
	\end{figure}
	Hình \ref{fig:activityuc06} mô tả các bước nông dân nhập thông tin thiết bị, chọn cảm biến liên kết và đặt điều kiện tự động hóa. Hệ thống kiểm tra quyền khu vực theo cảm biến và lưu thiết bị nếu hợp lệ.
	
	\paragraph{Biểu đồ tuần tự}
	\noindent
	\begin{figure}[H]
		\centering
		\safeincludegraphics[width=0.9\linewidth]{images/sequence_uc06}
		\caption{Biểu đồ tuần tự thêm thiết bị thông minh}
		\label{fig:sequenceuc06}
	\end{figure}
	Hình \ref{fig:sequenceuc06} mô tả quá trình giao diện gửi thông tin thiết bị đến máy chủ. Máy chủ kiểm tra quyền của nông dân đối với cảm biến liên kết, lưu cấu hình thiết bị và trả kết quả tạo mới cho giao diện.
	
	% -------------------------------------------------------
	\subsubsection{Ca sử dụng tự động kích hoạt thiết bị thông minh}
	
	\paragraph{Biểu đồ ca sử dụng}
	\noindent
	\begin{figure}[H]
		\centering
		\safeincludegraphics[width=0.85\linewidth]{images/uc_uc07}
		\caption{Biểu đồ ca sử dụng tự động kích hoạt thiết bị thông minh}
		\label{fig:ucuc07}
	\end{figure}
	Hình \ref{fig:ucuc07} mô tả ca sử dụng tự động kích hoạt thiết bị thông minh. Hệ thống nền thu thập dữ liệu cảm biến, so sánh với ngưỡng của cây trồng và bật thiết bị phù hợp khi phát hiện môi trường vượt hoặc thấp hơn ngưỡng.
	
	\begin{table}[H]
		\centering
		\caption{Đặc tả ca sử dụng tự động kích hoạt thiết bị thông minh}
		\label{tab:usecase-tudongkichhoat}
		\begin{tabular}{|p{4cm}|p{10cm}|}
			\hline
			\textbf{Tên ca sử dụng} & \textbf{Tự động kích hoạt thiết bị thông minh} \\
			\hline
			Tác nhân chính & Hệ thống nền \\
			\hline
			Dòng sự kiện chính &
			\begin{minipage}[t]{10cm}
				\vspace{2pt}
				\begin{enumerate}
					\item Hệ thống nền chạy định kỳ mỗi 60 giây.
					\item Hệ thống lưu số liệu mới của cảm biến.
					\item Hệ thống lấy ngưỡng trong hồ sơ cây trồng của khu vực.
					\item Hệ thống kiểm tra dữ liệu thấp hơn ngưỡng tối thiểu hoặc cao hơn ngưỡng tối đa.
					\item Nếu có vi phạm, hệ thống tạo hoặc cập nhật cảnh báo.
					\item Hệ thống tìm thiết bị đang bật tự động hóa, liên kết với đúng cảm biến và có điều kiện kích hoạt phù hợp.
					\item Hệ thống bật thiết bị, ghi kết quả tự động hóa vào cảnh báo và tự tắt thiết bị sau thời gian đã cấu hình.
				\end{enumerate}
				\vspace{2pt}
			\end{minipage} \\
			\hline
			Dòng sự kiện phụ & Nếu không có hồ sơ cây trồng, không có vi phạm ngưỡng hoặc không có thiết bị phù hợp, hệ thống chỉ lưu dữ liệu và không bật thiết bị. \\
			\hline
			Điều kiện tiên quyết & Khu vực có cảm biến hoạt động; thiết bị đã liên kết cảm biến và đã bật tự động hóa. \\
			\hline
			Điều kiện sau & Cảnh báo được lưu; thiết bị phù hợp được bật tự động và tự tắt sau thời gian đã cấu hình. \\
			\hline
		\end{tabular}
	\end{table}
	
	\paragraph{Biểu đồ hoạt động}
	\noindent
	\begin{figure}[H]
		\centering
		\safeincludegraphics[width=0.6\linewidth]{images/activity_uc07}
		\caption{Biểu đồ hoạt động tự động kích hoạt thiết bị thông minh}
		\label{fig:activityuc07}
	\end{figure}
	Hình \ref{fig:activityuc07} mô tả các bước hệ thống nền lưu dữ liệu cảm biến, kiểm tra ngưỡng và tìm thiết bị phù hợp. Nếu thiết bị có điều kiện kích hoạt khớp với hướng vi phạm, hệ thống bật thiết bị và tự tắt sau thời gian cấu hình.
	
	\paragraph{Biểu đồ tuần tự}
	\noindent
	\begin{figure}[H]
		\centering
		\safeincludegraphics[width=0.9\linewidth]{images/sequence_07}
		\caption{Biểu đồ tuần tự tự động kích hoạt thiết bị thông minh}
		\label{fig:sequence07}
	\end{figure}
	Hình \ref{fig:sequence07} mô tả quá trình bộ lập lịch gọi tác vụ nền để xử lý dữ liệu cảm biến. Tác vụ này lưu số liệu đo, đánh giá ngưỡng, ghi nhận cảnh báo, tìm thiết bị phù hợp với cảm biến liên kết, bật thiết bị khi cần can thiệp và tự tắt sau thời gian đã cấu hình.
	
	
	\newpage
	\subsection{Mô hình hóa dữ liệu}
	
	\paragraph{Sơ đồ thiết kế cơ sở dữ liệu}
	\noindent
	\begin{figure}[H]
		\centering
		\safeincludegraphics[width=0.95\linewidth]{images/DB}
		\caption{Mô hình cơ sở dữ liệu}
		\label{fig:db}
	\end{figure}
	
	\paragraph{Danh sách các đối tượng}
	Hệ thống AgriSmart sử dụng 8 bảng chính để lưu trữ thông tin người dùng, khu vực trồng trọt, cảm biến, thiết bị điều khiển và lịch sử cảnh báo môi trường. Danh sách các bảng được trình bày trong Bảng \ref{tab:doi-tuong}.
	
	\begin{table}[H]
		\centering
		\caption{Các đối tượng dữ liệu chính}
		\label{tab:doi-tuong}
		\begingroup
		\setlength{\tabcolsep}{4pt}
		\begin{tabularx}{\textwidth}{|>{\centering\arraybackslash}p{1cm}|>{\raggedright\arraybackslash}p{4cm}|Y|}
			\hline
			\centering\textbf{STT} & \centering\textbf{Tên bảng} & \centering\arraybackslash\textbf{Mô tả} \\
			\hline
			1 & users & Lưu thông tin tài khoản người dùng (quản trị viên và nông dân). \\
			\hline
			2 & growing\_zones & Lưu thông tin các khu vực trồng trọt. \\
			\hline
			3 & plant\_profiles & Lưu hồ sơ cây trồng và ngưỡng thông số môi trường tối ưu. \\
			\hline
			4 & farmer\_zone\_assignments & Bảng phân công nông dân quản lý khu vực (nhiều-nhiều). \\
			\hline
			5 & sensors & Lưu thông tin các cảm biến môi trường trong khu vực. \\
			\hline
			6 & sensor\_readings & Lưu bản ghi đo của cảm biến theo chu kỳ thời gian. \\
			\hline
			7 & devices & Lưu thông tin thiết bị thông minh, cảm biến liên kết và cấu hình tự động hóa. \\
			\hline
			8 & alerts & Lưu cảnh báo môi trường và trạng thái kích hoạt tự động hóa. \\
			\hline
		\end{tabularx}
		\endgroup
	\end{table}
	
	
	\begin{table}[H]
		\centering
		\caption{Bảng users}
		\label{tab:users}
		\begingroup
		\setlength{\tabcolsep}{3pt}
		\begin{tabularx}{\textwidth}{|>{\centering\arraybackslash}p{1cm}|>{\raggedright\arraybackslash}p{3.8cm}|>{\raggedright\arraybackslash}p{3.2cm}|Y|}
			\hline
			\centering\textbf{STT} & \centering\textbf{Tên thuộc tính} & \centering\textbf{Kiểu dữ liệu} & \centering\arraybackslash\textbf{Mô tả} \\
			\hline
			1 & id & UUID & Mã người dùng, khóa chính. \\
			\hline
			2 & email & VARCHAR(255) & Email đăng nhập, duy nhất và được đánh chỉ mục. \\
			\hline
			3 & hashed\_password & VARCHAR(255) & Mật khẩu đã được băm bằng bcrypt. \\
			\hline
			4 & full\_name & VARCHAR(255) & Họ và tên người dùng. \\
			\hline
			5 & role & ENUM & Vai trò: admin hoặc user (farmer). \\
			\hline
			6 & status & ENUM & Trạng thái tài khoản: active, inactive hoặc banned. \\
			\hline
			7 & is\_verified & BOOLEAN & Đã xác minh địa chỉ email hay chưa. \\
			\hline
			8 & created\_at & TIMESTAMPTZ & Thời điểm tạo tài khoản. \\
			\hline
			9 & updated\_at & TIMESTAMPTZ & Thời điểm cập nhật gần nhất. \\
			\hline
		\end{tabularx}
		\endgroup
	\end{table}
	
	\begin{table}[H]
		\centering
		\caption{Bảng growing\_zones}
		\label{tab:growing_zones}
		\begingroup
		\setlength{\tabcolsep}{3pt}
		\begin{tabularx}{\textwidth}{|>{\centering\arraybackslash}p{1cm}|>{\raggedright\arraybackslash}p{3.8cm}|>{\raggedright\arraybackslash}p{3.2cm}|Y|}
			\hline
			\centering\textbf{STT} & \centering\textbf{Tên thuộc tính} & \centering\textbf{Kiểu dữ liệu} & \centering\arraybackslash\textbf{Mô tả} \\
			\hline
			1 & id & UUID & Mã khu vực trồng trọt, khóa chính. \\
			\hline
			2 & name & VARCHAR(255) & Tên khu vực. \\
			\hline
			3 & description & VARCHAR(1000) & Mô tả khu vực. \\
			\hline
			4 & location & VARCHAR(255) & Vị trí địa lý. \\
			\hline
			5 & area\_sqm & FLOAT & Diện tích khu vực (m\textsuperscript{2}). \\
			\hline
			6 & plant\_profile\_id & UUID (FK) & Hồ sơ cây trồng đang được áp dụng. \\
			\hline
			7 & is\_active & BOOLEAN & Trạng thái hoạt động của khu vực. \\
			\hline
			8 & planting\_date & DATE & Ngày bắt đầu trồng. \\
			\hline
			9 & expected\_harvest\_date & DATE & Ngày thu hoạch dự kiến. \\
			\hline
			10 & created\_at / updated\_at & TIMESTAMPTZ & Thời điểm tạo và cập nhật bản ghi. \\
			\hline
		\end{tabularx}
		\endgroup
	\end{table}
	
	\begin{table}[H]
		\centering
		\caption{Bảng plant\_profiles}
		\label{tab:plant_profiles}
		\begingroup
		\setlength{\tabcolsep}{3pt}
		\begin{tabularx}{\textwidth}{|>{\centering\arraybackslash}p{1cm}|>{\raggedright\arraybackslash}p{3.8cm}|>{\raggedright\arraybackslash}p{3.2cm}|Y|}
			\hline
			\centering\textbf{STT} & \centering\textbf{Tên thuộc tính} & \centering\textbf{Kiểu dữ liệu} & \centering\arraybackslash\textbf{Mô tả} \\
			\hline
			1 & id & UUID & Mã hồ sơ cây trồng, khóa chính. \\
			\hline
			2 & name & VARCHAR(255) & Tên cây trồng, ràng buộc duy nhất. \\
			\hline
			3 & description & VARCHAR(1000) & Mô tả đặc điểm sinh trưởng. \\
			\hline
			4 & temp\_min / temp\_max & FLOAT & Dải nhiệt độ không khí tối ưu (\textdegree C). \\
			\hline
			5 & humidity\_min / max & FLOAT & Dải độ ẩm không khí tối ưu (\%). \\
			\hline
			6 & soil\_moisture\_min / max & FLOAT & Dải độ ẩm đất tối ưu (\%). \\
			\hline
			7 & light\_min / light\_max & FLOAT & Dải cường độ ánh sáng tối ưu (lux). \\
			\hline
			8 & ph\_min / ph\_max & FLOAT & Dải pH đất tối ưu. \\
			\hline
			9 & ec\_min / ec\_max & FLOAT & Dải độ dẫn điện tối ưu ($\mu$S/cm). \\
			\hline
			10 & growth\_period\_days & INTEGER & Thời gian sinh trưởng điển hình (ngày). \\
			\hline
			11 & created\_at / updated\_at & TIMESTAMPTZ & Thời điểm tạo và cập nhật hồ sơ. \\
			\hline
		\end{tabularx}
		\endgroup
	\end{table}
	
	\begin{table}[H]
		\centering
		\caption{Bảng farmer\_zone\_assignments}
		\label{tab:farmer_zone_assignments}
		\begingroup
		\setlength{\tabcolsep}{3pt}
		\begin{tabularx}{\textwidth}{|>{\centering\arraybackslash}p{1cm}|>{\raggedright\arraybackslash}p{3.8cm}|>{\raggedright\arraybackslash}p{3.2cm}|Y|}
			\hline
			\centering\textbf{STT} & \centering\textbf{Tên thuộc tính} & \centering\textbf{Kiểu dữ liệu} & \centering\arraybackslash\textbf{Mô tả} \\
			\hline
			1 & farmer\_id & UUID (PK, FK) & Khóa chính và khóa ngoại tới bảng users. \\
			\hline
			2 & zone\_id & UUID (PK, FK) & Khóa chính và khóa ngoại tới bảng growing\_zones. \\
			\hline
			3 & assigned\_at & TIMESTAMPTZ & Thời điểm phân công khu vực. \\
			\hline
		\end{tabularx}
		\endgroup
	\end{table}
	
	\begin{table}[H]
		\centering
		\caption{Bảng sensors}
		\label{tab:sensors}
		\begingroup
		\setlength{\tabcolsep}{3pt}
		\begin{tabularx}{\textwidth}{|>{\centering\arraybackslash}p{1cm}|>{\raggedright\arraybackslash}p{3.8cm}|>{\raggedright\arraybackslash}p{3.2cm}|Y|}
			\hline
			\centering\textbf{STT} & \centering\textbf{Tên thuộc tính} & \centering\textbf{Kiểu dữ liệu} & \centering\arraybackslash\textbf{Mô tả} \\
			\hline
			1 & id & UUID & Mã cảm biến, khóa chính. \\
			\hline
			2 & name & VARCHAR(255) & Tên cảm biến. \\
			\hline
			3 & sensor\_type & ENUM & Loại: temperature, humidity, soil\_moisture, light, pH, EC, CO2. \\
			\hline
			4 & unit & VARCHAR(50) & Đơn vị đo của cảm biến. \\
			\hline
			5 & zone\_id & UUID (FK) & Khóa ngoại tới khu vực trồng trọt. \\
			\hline
			6 & is\_active & BOOLEAN & Trạng thái hoạt động. \\
			\hline
			7 & created\_at / updated\_at & TIMESTAMPTZ & Thời điểm tạo và cập nhật. \\
			\hline
		\end{tabularx}
		\endgroup
	\end{table}
	
	\begin{table}[H]
		\centering
		\caption{Bảng sensor\_readings}
		\label{tab:sensor_readings}
		\begingroup
		\setlength{\tabcolsep}{3pt}
		\begin{tabularx}{\textwidth}{|>{\centering\arraybackslash}p{1cm}|>{\raggedright\arraybackslash}p{3.8cm}|>{\raggedright\arraybackslash}p{3.2cm}|Y|}
			\hline
			\centering\textbf{STT} & \centering\textbf{Tên thuộc tính} & \centering\textbf{Kiểu dữ liệu} & \centering\arraybackslash\textbf{Mô tả} \\
			\hline
			1 & id & UUID & Mã bản ghi đo, khóa chính. \\
			\hline
			2 & sensor\_id & UUID (FK) & Khóa ngoại tới cảm biến nguồn. \\
			\hline
			3 & value & FLOAT & Giá trị đo được (đơn vị theo loại cảm biến). \\
			\hline
			4 & recorded\_at & TIMESTAMPTZ & Thời điểm ghi đo, được đánh chỉ mục. \\
			\hline
		\end{tabularx}
		\endgroup
	\end{table}
	
	\begin{table}[H]
		\centering
		\caption{Bảng devices}
		\label{tab:devices}
		\begingroup
		\setlength{\tabcolsep}{3pt}
		\begin{tabularx}{\textwidth}{|>{\centering\arraybackslash}p{1cm}|>{\raggedright\arraybackslash}p{3.8cm}|>{\raggedright\arraybackslash}p{3.2cm}|Y|}
			\hline
			\centering\textbf{STT} & \centering\textbf{Tên thuộc tính} & \centering\textbf{Kiểu dữ liệu} & \centering\arraybackslash\textbf{Mô tả} \\
			\hline
			1 & id & UUID & Mã thiết bị thông minh, khóa chính. \\
			\hline
			2 & name / location & VARCHAR & Tên thiết bị và vị trí lắp đặt. \\
			\hline
			3 & type & VARCHAR(50) & Loại thiết bị: pump, fan, heater, light, fertilizer pump, shade. \\
			\hline
			4 & owner\_id & UUID (FK) & Người sở hữu hoặc quản lý thiết bị. \\
			\hline
			5 & linked\_sensor\_id & UUID (FK) & Cảm biến dùng để kích hoạt tự động hóa. \\
			\hline
			6 & automation\_enabled & BOOLEAN & Bật hoặc tắt chế độ tự động hóa. \\
			\hline
			7 & automation\_trigger & VARCHAR(30) & Điều kiện kích hoạt: vượt ngưỡng, dưới ngưỡng hoặc cả hai. \\
			\hline
			8 & timeout\_seconds & INTEGER & Thời gian chạy trước khi thiết bị tự tắt. \\
			\hline
			9 & current\_state / current\_value & VARCHAR / FLOAT & Trạng thái hiện tại và giá trị điều khiển hiện tại. \\
			\hline
			10 & connection\_status & VARCHAR(30) & Trạng thái kết nối: online hoặc offline. \\
			\hline
			11 & last\_command & VARCHAR(500) & Lệnh gần nhất, gồm lệnh thủ công hoặc tự động. \\
			\hline
			12 & is\_active & BOOLEAN & Thiết bị còn được sử dụng trong hệ thống hay không. \\
			\hline
			13 & created\_at / updated\_at & TIMESTAMPTZ & Thời điểm tạo và cập nhật. \\
			\hline
		\end{tabularx}
		\endgroup
	\end{table}
	
	\begin{table}[H]
		\centering
		\caption{Bảng alerts}
		\label{tab:alerts}
		\begingroup
		\setlength{\tabcolsep}{3pt}
		\begin{tabularx}{\textwidth}{|>{\centering\arraybackslash}p{1cm}|>{\raggedright\arraybackslash}p{3.8cm}|>{\raggedright\arraybackslash}p{3.2cm}|Y|}
			\hline
			\centering\textbf{STT} & \centering\textbf{Tên thuộc tính} & \centering\textbf{Kiểu dữ liệu} & \centering\arraybackslash\textbf{Mô tả} \\
			\hline
			1 & id & UUID & Mã cảnh báo, khóa chính. \\
			\hline
			2 & zone\_id & UUID (FK) & Khu vực phát sinh cảnh báo. \\
			\hline
			3 & sensor\_id & UUID (FK) & Cảm biến liên quan, có thể rỗng. \\
			\hline
			4 & alert\_type & ENUM & Loại: above\_max, below\_min, compound\_condition, device\_offline. \\
			\hline
			5 & severity & ENUM & Mức độ nghiêm trọng: low, medium, high. \\
			\hline
			6 & parameter & ENUM & Thông số môi trường lệch ngưỡng (loại cảm biến). \\
			\hline
			7 & actual\_value & FLOAT & Giá trị thực tế đo được tại thời điểm kích hoạt. \\
			\hline
			8 & threshold\_value & FLOAT & Ngưỡng tham chiếu bị vi phạm. \\
			\hline
			9 & message & VARCHAR(500) & Thông điệp mô tả cảnh báo. \\
			\hline
			10 & recommended\_action & VARCHAR(500) & Hành động được khuyến nghị cho người vận hành. \\
			\hline
			11 & automation\_status & VARCHAR(30) & Trạng thái tự động hóa: none, executed hoặc skipped. \\
			\hline
			12 & automation\_device\_id & UUID (FK) & Thiết bị đã được kích hoạt tự động, nếu có. \\
			\hline
			13 & automation\_device\_name & VARCHAR(255) & Tên thiết bị đã được kích hoạt tự động. \\
			\hline
			14 & automation\_command & VARCHAR(255) & Lệnh tự động đã gửi cho thiết bị. \\
			\hline
			15 & is\_read & BOOLEAN & Đã được đọc và xử lý hay chưa. \\
			\hline
			16 & triggered\_at & TIMESTAMPTZ & Thời điểm kích hoạt cảnh báo, được đánh chỉ mục. \\
			\hline
		\end{tabularx}
		\endgroup
	\end{table}
	
	\chaptersection{XÂY DỰNG CHƯƠNG TRÌNH}
	
	\subsection{Kiến trúc hệ thống}
	
	AgriSmart được xây dựng theo mô hình kiến trúc \textbf{máy khách -- máy chủ ba tầng}, trong đó mỗi tầng có trách nhiệm riêng biệt và giao tiếp với nhau qua giao thức được định nghĩa rõ ràng, tạo nền tảng để mở rộng và bảo trì độc lập từng thành phần.
	
	\begin{figure}[H]
		\centering
		\begin{tikzpicture}[
			every node/.style={font=\normalsize},
			tbox/.style={draw=myblue, fill=lightblue, rounded corners=6pt,
				minimum width=3.8cm, minimum height=1cm, align=center,
				text=black},
			arrow/.style={-{Stealth[length=6pt]}, myblue, thick},
			darrow/.style={-{Stealth[length=6pt]}, myblue, thick, dashed}
			]
			% Ba hộp chính
			\node[tbox] (client) at (0,0)    {\textbf{Trình duyệt}\\(Tầng giao diện)};
			\node[tbox] (server) at (5.5,0)  {\textbf{Máy chủ}\\(Tầng xử lý)};
			\node[tbox] (db)     at (11,0)   {\textbf{Cơ sở dữ liệu}\\(Tầng lưu trữ)};
			
			% Mũi tên Client ↔ Server
			\draw[arrow]  ([yshift= 3pt]client.east) -- node[above, font=\normalsize]{Yêu cầu} ([yshift= 3pt]server.west);
			\draw[arrow]  ([yshift=-3pt]server.west) -- node[below, font=\normalsize]{Phản hồi}       ([yshift=-3pt]client.east);
			
			% Mũi tên Server ↔ DB
			\draw[arrow]  ([yshift= 3pt]server.east) -- node[above, font=\normalsize]{Truy vấn}   ([yshift= 3pt]db.west);
			\draw[arrow]  ([yshift=-3pt]db.west)     -- node[below, font=\normalsize]{Kết quả dữ liệu}       ([yshift=-3pt]server.east);
			
			% Tác vụ nền tự vòng lại
			\draw[darrow] (server.north) arc[start angle=0, end angle=180, radius=0.9cm]
			node[above, pos=0.5, font=\normalsize]{Tác vụ nền (60 giây)};
			
			% Nhãn tầng phía trên
			\node[above=1.6cm of client, font=\bfseries\normalsize, text=myblue] {TẦNG GIAO DIỆN};
			\node[above=1.6cm of server, font=\bfseries\normalsize, text=myblue] {TẦNG XỬ LÝ};
			\node[above=1.6cm of db,     font=\bfseries\normalsize, text=myblue] {TẦNG DỮ LIỆU};
			
			% Đường kẻ ngăn cách tầng (nét đứt nhạt)
			\draw[linegray, dashed, thin] (2.75, -0.9) -- (2.75, 1.0);
			\draw[linegray, dashed, thin] (8.25, -0.9) -- (8.25, 1.0);
		\end{tikzpicture}
		\caption{Kiến trúc ba tầng của hệ thống AgriSmart}
		\label{fig:kien-truc}
	\end{figure}
	
	Tầng \textbf{giao diện} là ứng dụng web chạy trên trình duyệt, nơi người dùng --- quản trị viên và nông dân --- thực hiện các thao tác nghiệp vụ như xem dữ liệu cảm biến, xử lý cảnh báo và điều khiển thiết bị. Mọi yêu cầu được gửi lên tầng xử lý qua giao thức HTTPS/REST và nhận lại dữ liệu dạng JSON để hiển thị.
	
	Tầng \textbf{xử lý} tiếp nhận yêu cầu, xác thực danh tính, kiểm tra quyền truy cập theo vai trò, áp dụng logic nghiệp vụ và điều phối tương tác với cơ sở dữ liệu. Ngoài việc phục vụ các yêu cầu từ giao diện, tầng này còn chạy một tác vụ nền định kỳ mỗi 60 giây để thu thập dữ liệu cảm biến, đánh giá bộ quy tắc sinh cảnh báo và kích hoạt thiết bị phù hợp một cách tự động.
	
	Tầng \textbf{dữ liệu} lưu trữ toàn bộ thông tin vận hành của hệ thống: tài khoản người dùng, khu vực trồng trọt, hồ sơ cây trồng, chuỗi số liệu cảm biến, thiết bị thông minh, cấu hình tự động hóa và cảnh báo môi trường. Tầng này chỉ được truy cập bởi tầng xử lý, không bao giờ bị truy cập trực tiếp từ giao diện.
	
	\subsection{Công nghệ sử dụng}
	
	\subsubsection{Python và FastAPI}
	
	Phía máy chủ, hệ thống sử dụng Python kết hợp FastAPI để xây dựng lớp API phục vụ các chức năng giám sát cảm biến, quản lý thiết bị thông minh và sinh cảnh báo môi trường. FastAPI được lựa chọn nhờ khả năng xử lý bất đồng bộ, phù hợp với mô hình tác vụ nền định kỳ thu thập và lưu dữ liệu cảm biến mỗi 60 giây, đồng thời đánh giá luật cảnh báo theo ngưỡng hồ sơ cây trồng và kích hoạt thiết bị phù hợp. FastAPI tích hợp Pydantic để kiểm tra kiểu dữ liệu đầu vào tại các chức năng như tạo hồ sơ cây trồng, cấu hình thiết bị và truy vấn lịch sử đo cảm biến, đồng thời hỗ trợ tách biệt rõ ràng các tầng xác thực, phân quyền và xử lý nghiệp vụ. Cơ chế xác thực và băm mật khẩu đảm bảo an toàn truy cập cho cả hai vai trò quản trị viên và nông dân.
	
	\begin{figure}[H]
		\centering
		\safeincludegraphics[width=0.75\linewidth]{images/fastapi}
		\caption{FastAPI -- Công cụ xây dựng API phía máy chủ}
		\label{fig:fastapi}
	\end{figure}
	
	\subsubsection{PostgreSQL}
	
	PostgreSQL được chọn làm hệ quản trị cơ sở dữ liệu quan hệ nhờ hỗ trợ đầy đủ các kiểu dữ liệu cần thiết và khả năng xử lý giao dịch nguyên tử -- điều quan trọng khi ghi đồng thời số liệu cảm biến, sinh cảnh báo và cập nhật trạng thái thiết bị trong cùng một chu kỳ nền. Toàn bộ tương tác được thực hiện qua SQLAlchemy ở chế độ bất đồng bộ, giúp các lớp xử lý dữ liệu dễ bảo trì và kiểm thử độc lập. Alembic quản lý lịch sử thay đổi lược đồ dữ liệu từ thiết kế ban đầu đến tái cấu trúc quan hệ phân công nông dân-khu vực, bổ sung loại cảnh báo tổ hợp và cấu hình tự động hóa thiết bị.
	
	\begin{figure}[H]
		\centering
		\safeincludegraphics[width=0.75\linewidth]{images/postgres}
		\caption{PostgreSQL -- Hệ quản trị cơ sở dữ liệu quan hệ}
		\label{fig:tech-postgresql}
	\end{figure}
	
	\subsubsection{React, TypeScript và Vite}
	
	Giao diện người dùng được xây dựng dưới dạng ứng dụng đơn trang sử dụng React 19 kết hợp TypeScript. React cho phép tổ chức màn hình theo hướng thành phần độc lập: thẻ tóm tắt trạng thái khu vực, biểu đồ xu hướng cảm biến theo chuỗi thời gian, bảng cảnh báo phân loại theo mức độ nghiêm trọng và bảng quản lý thiết bị thông minh đều có thể tái sử dụng trên nhiều khu vực trồng. TypeScript bổ sung hệ thống kiểu tĩnh, giúp phát hiện sớm sai lệch kiểu dữ liệu khi trao đổi với máy chủ. Vite cung cấp môi trường phát triển có khả năng tải lại nhanh, rút ngắn chu kỳ kiểm tra giao diện trong quá trình phát triển tính năng mới.
	
	\begin{figure}[H]
		\centering
		\safeincludegraphics[width=0.75\linewidth]{images/react-vite}
		\caption{React và TypeScript -- Nền tảng xây dựng giao diện người dùng}
		\label{fig:react-vite}
	\end{figure}
	
	\subsubsection{TanStack Router, TanStack Query và Zustand}
	
	TanStack Router đảm nhận điều hướng phía giao diện với cơ chế bảo vệ theo vai trò, đảm bảo quản trị viên mới truy cập được các trang quản trị (quản lý người dùng, hồ sơ cây trồng, phân công khu vực) trong khi nông dân chỉ thấy các khu vực được phân công cho tài khoản mình. TanStack Query quản lý vòng đời yêu cầu dữ liệu, tự động làm mới khi cần và hỗ trợ cập nhật giao diện nhanh -- đặc biệt phù hợp khi đồng bộ dữ liệu cảm biến cập nhật mỗi 60 giây và trạng thái cảnh báo thay đổi theo chu kỳ nền. Zustand lưu trữ trạng thái dùng chung của giao diện gồm thông tin phiên đăng nhập và danh sách cảnh báo chưa đọc.
	
	\subsubsection{Shadcn/UI}
	
	\begin{figure}[H]
		\centering
		\safeincludegraphics[width=0.75\linewidth]{images/shadcn}
		\caption{Shadcn UI}
		\label{fig:shadcn}
	\end{figure}
	
	Shadcn/UI cung cấp bộ thành phần giao diện sẵn có như hộp thoại, danh sách chọn, nhãn trạng thái, bảng dữ liệu, biểu mẫu và thông báo nhanh, phục vụ các màn hình như tạo khu vực trồng, bảng cảnh báo phân loại theo mức cao/trung bình/thấp và cấu hình thiết bị thông minh. Tailwind CSS định dạng toàn bộ giao diện theo hướng tiện ích, đảm bảo nhất quán về màu sắc phân loại cảnh báo (đỏ-cam-vàng) và trạng thái thiết bị (xanh-xám) trên mọi màn hình mà không cần tệp CSS riêng. Recharts hiển thị biểu đồ xu hướng các thông số cảm biến (nhiệt độ, độ ẩm, ánh sáng, pH, EC, CO\textsubscript{2}) theo chuỗi thời gian, giúp người vận hành dễ dàng nhận ra xu hướng bất thường trong ngày canh tác.
	
	
	\subsubsection{Docker}
	
	Toàn bộ hệ thống được đóng gói và vận hành bằng Docker Compose với năm dịch vụ: phần máy chủ dùng FastAPI kèm tác vụ nền mô phỏng cảm biến, phần giao diện dùng React/Nginx, cơ sở dữ liệu PostgreSQL, Redis phục vụ lưu tạm và MailHog giả lập gửi email đặt lại mật khẩu trong môi trường phát triển. Mỗi dịch vụ chạy trong vùng chứa riêng với biến môi trường được quản lý qua tệp cấu hình, đảm bảo mọi môi trường phát triển và triển khai đều hoạt động trên cùng cấu hình. Cách tiếp cận này cho phép khởi chạy toàn bộ hệ thống với một lệnh duy nhất và tạo nền thuận lợi để tích hợp thêm cổng kết nối IoT khi mở rộng sang phần cứng thực địa.
	
	\begin{figure}[H]
		\centering
		\safeincludegraphics[width=0.75\linewidth]{images/docker}
		\caption{Docker Compose -- Môi trường container hóa và triển khai hệ thống}
		\label{fig:docker}
	\end{figure}
	
	
	\subsection{Giao diện hệ thống}
	
	Các giao diện của hệ thống được trình bày theo đúng thứ tự các ca sử dụng chính: đăng nhập, tạo hồ sơ cây trồng, tạo khu vực trồng trọt, phân công nông dân vào khu vực, thêm cảm biến, thêm thiết bị thông minh, sau đó theo dõi cảnh báo và kết quả tự động hóa. Cách sắp xếp này giúp người đọc nhìn được luồng thiết lập và vận hành hệ thống từ bước đầu tiên đến khi dữ liệu cảm biến được dùng để kích hoạt thiết bị.
	
	\subsubsection{Giao diện đăng nhập}
	Hình \ref{fig:uidangnhap} mô tả màn hình đăng nhập của hệ thống. Người dùng nhập email và mật khẩu đã được cấp; nếu thông tin hợp lệ, hệ thống xác định vai trò của người dùng và điều hướng đến màn hình phù hợp. Nếu thông tin sai hoặc tài khoản không còn hoạt động, hệ thống hiển thị thông báo lỗi để người dùng kiểm tra lại.
	
	\begin{figure}[H]
		\centering
		\safeincludegraphics[width=1\linewidth]{images/ui_dang_nhap}
		\caption{Giao diện đăng nhập}
		\label{fig:uidangnhap}
	\end{figure}
	
	\subsubsection{Giao diện tạo hồ sơ cây trồng}
	Hình \ref{fig:uihscaytrong} mô tả màn hình tạo và quản lý hồ sơ cây trồng. Quản trị viên nhập tên cây trồng, mô tả và các dải ngưỡng môi trường như nhiệt độ, độ ẩm, độ ẩm đất, ánh sáng, pH, độ dẫn điện và CO\textsubscript{2}. Hồ sơ cây trồng là cơ sở để hệ thống đánh giá dữ liệu cảm biến và sinh cảnh báo ở các bước vận hành sau.
	
	\begin{figure}[H]
		\centering
		\safeincludegraphics[width=1\linewidth]{images/ui_hs_cay_trong}
		\caption{Giao diện hồ sơ cây trồng với danh sách hồ sơ và biểu mẫu nhập các dải ngưỡng môi trường tối ưu}
		\label{fig:uihscaytrong}
	\end{figure}
	
	\subsubsection{Giao diện tạo khu vực trồng trọt}
	Sau khi có hồ sơ cây trồng, quản trị viên tạo khu vực trồng trọt để xác định phạm vi theo dõi thực tế. Hình \ref{fig:uitaokvtrongtrot} mô tả biểu mẫu tạo khu vực với các thông tin như tên khu vực, mô tả, vị trí, diện tích, ngày bắt đầu trồng và hồ sơ cây trồng áp dụng. Việc gán hồ sơ cây trồng ngay khi tạo khu vực giúp hệ thống biết ngưỡng môi trường nào cần dùng khi đánh giá dữ liệu cảm biến của khu vực đó.
	
	\begin{figure}[H]
		\centering
		\safeincludegraphics[width=0.7\linewidth]{images/ui_tao_kv_trong_trot}
		\caption{Giao diện tạo khu vực trồng trọt}
		\label{fig:uitaokvtrongtrot}
	\end{figure}
	
	\subsubsection{Giao diện phân công nông dân vào khu vực}
	Hình \ref{fig:uikhuvuctrongtrot} mô tả màn hình danh sách khu vực trồng trọt và phần phân công nông dân phụ trách từng khu vực. Sau khi khu vực được tạo, quản trị viên chọn nông dân phù hợp để cấp quyền xem dữ liệu, thêm cảm biến và quản lý thiết bị trong khu vực đó. Nhờ bước phân công này, nông dân chỉ thao tác trên đúng các khu vực được giao, giúp dữ liệu và thiết bị của từng vùng trồng được quản lý rõ ràng hơn.
	
	\begin{figure}[H]
		\centering
		\safeincludegraphics[width=1\linewidth]{images/ui_khu_vuc_trong_trot}
		\caption{Giao diện phân công nông dân vào khu vực}
		\label{fig:uikhuvuctrongtrot}
	\end{figure}
	
	\subsubsection{Giao diện thêm cảm biến}
	
	\begin{figure}[H]
		\centering
		\safeincludegraphics[width=0.7\linewidth]{images/ui_them_cam_bien}
		\caption{Giao diện thêm cảm biến}
		\label{fig:uithemcambien}
	\end{figure}
	Hình \ref{fig:uithemcambien} mô tả màn hình thêm một cảm biến cho khu vực trồng trọt. Người dùng chọn khu vực, nhập tên cảm biến, loại cảm biến, đơn vị đo và vị trí lắp đặt. Sau khi cảm biến được thêm, hệ thống có thể thu thập dữ liệu theo chu kỳ và hiển thị lịch sử đo dưới dạng biểu đồ để người vận hành theo dõi biến động môi trường.
	
	\begin{figure}[H]
		\centering
		\safeincludegraphics[width=1\linewidth]{images/ui_ds_cam_bien}
		\caption{Giao diện danh sách cảm biến}
		\label{fig:uidscambien}
	\end{figure}
	Hình \ref{fig:uidscambien} mô tả danh sách các cảm biến đã được lắp đặt trong hệ thống. Mỗi cảm biến hiển thị tên, loại thông số đo, khu vực liên kết, trạng thái hoạt động và thời điểm cập nhật gần nhất. Từ danh sách này, người vận hành có thể kiểm tra cảm biến nào đang phục vụ khu vực nào trước khi xem lịch sử đo hoặc liên kết cảm biến với thiết bị thông minh.
	
	\subsubsection{Giao diện thêm thiết bị thông minh}
	\begin{figure}[H]
		\centering
		\safeincludegraphics[width=0.7\linewidth]{images/ui_them_thiet_bi}
		\caption{Giao diện thêm thiết bị thông minh}
		\label{fig:uithemthietbi}
	\end{figure}
	Hình \ref{fig:uithemthietbi} mô tả biểu mẫu thêm thiết bị thông minh cho khu vực trồng trọt. Khi khai báo thiết bị, nông dân nhập tên thiết bị, loại thiết bị, vị trí lắp đặt, cảm biến liên kết và trạng thái tự động hóa. Nếu bật tự động hóa, người dùng tiếp tục chọn điều kiện kích hoạt như dưới ngưỡng, vượt ngưỡng hoặc cả hai, đồng thời cấu hình thời gian chạy để thiết bị tự tắt sau khi can thiệp.
	
	\begin{figure}[H]
		\centering
		\safeincludegraphics[width=1\linewidth]{images/ui_dk_tb_chap_hanh}
		\caption{Giao diện thiết bị thông minh với bộ lọc khu vực, cấu hình tự động hóa và điều khiển thủ công}
		\label{fig:uidktbchaphanh}
	\end{figure}
	
	Hình \ref{fig:uidktbchaphanh} mô tả màn hình quản lý các thiết bị thông minh sau khi đã được thêm vào hệ thống. Danh sách thiết bị có thể lọc theo khu vực trồng trọt, hiển thị cảm biến liên kết, trạng thái hiện tại và chế độ tự động hóa. Người vận hành có thể bật hoặc tắt thiết bị thủ công khi cần, đồng thời kiểm tra thiết bị nào đang được cấu hình để tự động phản ứng với dữ liệu cảm biến.
	
	\subsubsection{Giao diện cảnh báo và tự động kích hoạt thiết bị}
	Hình \ref{fig:uicanhbao} mô tả màn hình cảnh báo môi trường được sinh ra từ dữ liệu cảm biến. Mỗi cảnh báo hiển thị khu vực, thông số vi phạm, giá trị đo thực tế, mức độ nghiêm trọng, thời điểm phát sinh và khuyến nghị xử lý. Khi cảnh báo phù hợp với cấu hình thiết bị ở Hình \ref{fig:uidktbchaphanh}, hệ thống có thể tự động bật thiết bị tương ứng, ghi nhận thiết bị đã được kích hoạt và cho phép người vận hành đánh dấu cảnh báo đã đọc sau khi kiểm tra.
	
	\begin{figure}[H]
		\centering
		\safeincludegraphics[width=1\linewidth]{images/ui_canh_bao}
		\caption{Giao diện cảnh báo và tự động kích hoạt thiết bị}
		\label{fig:uicanhbao}
	\end{figure}
	
	\subsubsection{Giao diện tổng quan vận hành}
	Sau khi các khu vực, cảm biến và thiết bị đã được thiết lập, người dùng theo dõi trạng thái vận hành trên màn hình tổng quan. Hình \ref{fig:uitongquanadmin} mô tả giao diện tổng quan với các chỉ số chính như số khu vực, số cảm biến, số thiết bị, số cảnh báo chưa đọc và các biểu đồ số liệu mới nhất. Màn hình này giúp người vận hành nhanh chóng nắm được tình trạng chung của toàn hệ thống sau khi các bước trong luồng use case đã hoàn thành.
	
	\begin{figure}[H]
		\centering
		\safeincludegraphics[width=1\linewidth]{images/ui_tong_quan_admin}
		\caption{Giao diện trang tổng quan của quản trị viên}
		\label{fig:uitongquanadmin}
	\end{figure}
	
	\newpage
	\begin{center}
		\section*{KẾT LUẬN}
	\end{center}
	\addcontentsline{toc}{section}{KẾT LUẬN}
	
	Qua quá trình nghiên cứu, phân tích và phát triển, đồ án đã đạt được mục tiêu xây dựng hệ thống giám sát nông nghiệp thông minh trên nền tảng web. Hệ thống được hình thành từ nhu cầu thực tế trong quá trình theo dõi môi trường canh tác tại các mô hình nhà lưới, nhà kính và nông hộ quy mô vừa và nhỏ, hướng tới việc thay thế dần cách ghi chép thủ công, quan sát rời rạc và điều khiển thiết bị dựa hoàn toàn vào kinh nghiệm cá nhân. Trên cơ sở khảo sát nghiệp vụ và tham khảo các tài liệu khoa học về nông nghiệp chính xác, đồ án đã xác định và triển khai các chức năng cần thiết như đăng nhập, tạo hồ sơ cây trồng, tạo khu vực trồng trọt, phân công nông dân vào khu vực, thêm cảm biến, thêm thiết bị thông minh và tự động kích hoạt thiết bị khi môi trường vượt hoặc thấp hơn ngưỡng.
	
	Kết quả chính của đồ án là xây dựng được một hệ thống có thể phục vụ đồng thời hai nhóm người dùng: quản trị viên và nông dân. Quản trị viên có thể thiết lập hồ sơ cây trồng với các ngưỡng môi trường phù hợp, tạo khu vực trồng trọt và phân công nông dân phụ trách từng khu vực. Nông dân có thể theo dõi các khu vực được giao, thêm cảm biến, thêm thiết bị thông minh, liên kết thiết bị với cảm biến và cấu hình điều kiện tự động hóa. Khi dữ liệu cảm biến cho thấy môi trường vượt ngưỡng, hệ thống có thể sinh cảnh báo và kích hoạt thiết bị phù hợp, chẳng hạn bật quạt khi nhiệt độ cao hoặc bật máy sưởi khi nhiệt độ thấp.
	
	Bên cạnh việc hoàn thiện các chức năng cơ bản, hệ thống còn góp phần khắc phục một số hạn chế trong quy trình giám sát nông nghiệp thủ công. Dữ liệu cảm biến được tổ chức tập trung, các ngưỡng sinh trưởng được quản lý theo từng hồ sơ cây trồng, quyền thao tác của nông dân được kiểm soát theo khu vực và việc kích hoạt thiết bị được thực hiện dựa trên quy tắc rõ ràng. Nhờ đó, hệ thống hỗ trợ giảm độ trễ trong phát hiện bất thường, giúp người vận hành có căn cứ can thiệp kịp thời, đồng thời tạo nền tảng để mở rộng sang mô hình giám sát và điều khiển nông nghiệp chính xác trong tương lai.
	
	Tuy nhiên, do giới hạn về thời gian và phạm vi thực hiện, hệ thống vẫn còn một số điểm cần tiếp tục hoàn thiện. Dữ liệu cảm biến hiện mới được mô phỏng bằng phần mềm, chưa đọc trực tiếp từ thiết bị đo thực địa; cơ chế tự động hóa mới dựa trên ngưỡng và thời gian chạy cấu hình sẵn, chưa có phản hồi khép kín để kiểm tra môi trường đã trở về trạng thái an toàn sau khi bật thiết bị; giao diện vẫn cần tiếp tục tối ưu cho điện thoại vì nông dân thường thao tác trực tiếp tại khu vực canh tác. Những hạn chế này là cơ sở để hệ thống tiếp tục được phát triển theo hướng tích hợp thiết bị IoT thực, bổ sung cơ chế phản hồi sau điều khiển, mở rộng bộ quy tắc theo mùa vụ và giai đoạn sinh trưởng, cải thiện trải nghiệm trên thiết bị di động và thử nghiệm trong môi trường canh tác thực tế để thu thập phản hồi, từ đó điều chỉnh hệ thống phù hợp hơn với nhu cầu sử dụng.
	
	\newpage
	\phantomsection
	\addcontentsline{toc}{section}{DANH MỤC TÀI LIỆU THAM KHẢO}
	\begin{thebibliography}{99}
		
		\bibitem{boKhoaHoc2025a}
		Bộ Khoa học và Công nghệ (2025). \textit{Tầm nhìn chiến lược của Nghị quyết 57 cho nông nghiệp Việt Nam: Sơn La đột phá trong phát triển nông nghiệp công nghệ cao}. Truy cập tại: \url{https://mst.gov.vn/tam-nhin-chien-luoc-cua-nghi-quyet-57-cho-nong-nghiep-viet-nam-son-la-dot-pha.htm}
		
		\bibitem{soussi2024}
		Soussi, A., Zero, E., Sacile, R., Trinchero, D. và Fossa, M. (2024). Smart Sensors and Smart Data for Precision Agriculture: A Review. \textit{Sensors}.
		
		\bibitem{koksal2019}
		Koksal, O. và Tekinerdogan, B. (2019). Architecture Design Approach for IoT-Based Farm Management Information Systems. \textit{Precision Agriculture}, 20, 926--958.
		
		\bibitem{ojha2015}
		Ojha, T., Misra, S. và Raghuwanshi, N.S. (2015). Wireless Sensor Networks for Agriculture: The State-of-the-Art in Practice and Future Challenges. \textit{Computers and Electronics in Agriculture}, 118, 66--84.
		
		\bibitem{wolfert2017}
		Wolfert, S., Ge, L., Verdouw, C. và Bogaardt, M.J. (2017). Big Data in Smart Farming -- A Review. \textit{Agricultural Systems}, 153, 69--80.
		
		\bibitem{goap2018}
		Goap, A., Sharma, D., Shukla, A.K. và Kumar, C.R. (2018). An IoT Based Smart Irrigation Management System Using Machine Learning and Open Source Technologies. \textit{Computers and Electronics in Agriculture}, 155, 41--49.
		
		\bibitem{cucThuongMai2018}
		Cục Thương mại Biên giới và Miền núi (2018). \textit{Sơn La đã phát triển nhiều sản phẩm nông nghiệp mang tính đặc trưng của vùng miền}. Truy cập tại: \url{https://thuongmaibiengioimiennui.gov.vn/san-pham-vung-mien/mien-bac/2018/8/son-la-da-phat-trien-nhieu-san-pham-nong-nghiep-mang-tinh-dac-trung-cua-vung-mien}
		
		\bibitem{boKhoaHoc2025b}
		Bộ Khoa học và Công nghệ (2025). \textit{Nhận diện rào cản để người dân vùng cao tiếp cận chuyển đổi số}. Truy cập tại: \url{https://mst.gov.vn/nhan-dien-rao-can-de-nguoi-dan-vung-cao-tiep-can-chuyen-doi-so-197251109191.htm}
		
		\bibitem{ahmed2018}
		Ahmed, N., De, D. và Hussain, I. (2018). Internet of Things (IoT) for Smart Precision Agriculture and Farming in Rural Areas. \textit{IEEE Internet of Things Journal}, 5(6), 4890--4899.
		
		\bibitem{growlink2024}
		Growlink (2024). \textit{Smart Farm Automation and Monitoring Platform}. Truy cập tại: \url{https://www.growlink.ag}
		
		\bibitem{farmapp2024}
		Farmapp (2024). \textit{Farmapp -- Smart Agriculture IoT Platform}. Truy cập tại: \url{https://farmappweb.com/}
		
	\end{thebibliography}
	
	
	\newpage
	\phantomsection
	\addcontentsline{toc}{section}{PHỤ LỤC}
	\begin{center}
		\fontsize{14pt}{20pt}\selectfont
		\textbf{PHỤ LỤC}
	\end{center}
	
	\noindent\textbf{Nội dung gồm:}
	
	\noindent
	\begin{tabular}{@{}l@{}}
		Phụ lục A: Danh sách API hệ thống\\
		Phụ lục B: Script cơ sở dữ liệu (PostgreSQL)\\
		Phụ lục C: Code minh họa Backend (FastAPI)\\
		Phụ lục D: Frontend (React + TypeScript)\\
		Phụ lục E: Bảo mật hệ thống\\
		Phụ lục F: Triển khai hệ thống
	\end{tabular}
	
	\phantomsection
	\addcontentsline{toc}{subsection}{Phụ lục A: Danh sách API hệ thống}
	\subsection*{Phụ lục A: Danh sách API hệ thống}
	
	\subsubsection*{A.1. Quy ước chung}
	
	\noindent
	Base URL: \texttt{/api/v1}
	
	\noindent
	Dữ liệu trao đổi: JSON
	
	\noindent
	Xác thực: JWT lưu trong HTTP-only cookie; khi gọi API bằng công cụ kiểm thử có thể truyền theo Bearer Token.
	
	\noindent
	Header:
	\begin{itemize}
		\item \texttt{Authorization: Bearer \textless token\textgreater}
		\item \texttt{Content-Type: application/json}
	\end{itemize}
	
	\subsubsection*{A.2. Response chuẩn}
	
	\noindent
	\textbf{Thành công}
	
	\noindent
	\begin{lstlisting}
		{ "success": true, "message": "Thanh cong", "data": {} }
	\end{lstlisting}
	
	\noindent
	\textbf{Thất bại}
	
	\noindent
	\begin{lstlisting}
		{ "success": false, "message": "Noi dung loi", "data": null }
	\end{lstlisting}
	
	\subsubsection*{A.3. Các nhóm API chính}
	
	\noindent
	\textbf{Xác thực và tài khoản}
	\begin{itemize}
		\item \texttt{POST /auth/register} -- đăng ký tài khoản
		\item \texttt{POST /auth/login} -- đăng nhập
		\item \texttt{POST /auth/logout} -- đăng xuất
		\item \texttt{GET /auth/me} -- lấy thông tin người dùng hiện tại
		\item \texttt{PUT /auth/change-password} -- đổi mật khẩu
	\end{itemize}
	
	\noindent
	\textbf{Hồ sơ cây trồng}
	\begin{itemize}
		\item \texttt{GET /plant-profiles} -- danh sách hồ sơ cây trồng
		\item \texttt{GET /plant-profiles/\{profile\_id\}} -- chi tiết hồ sơ cây trồng
		\item \texttt{POST /plant-profiles} -- tạo hồ sơ cây trồng
		\item \texttt{PUT /plant-profiles/\{profile\_id\}} -- cập nhật hồ sơ cây trồng
		\item \texttt{DELETE /plant-profiles/\{profile\_id\}} -- xóa hồ sơ cây trồng
	\end{itemize}
	
	\noindent
	\textbf{Khu vực trồng trọt}
	\begin{itemize}
		\item \texttt{GET /zones} -- danh sách khu vực được phân công
		\item \texttt{GET /zones/\{zone\_id\}} -- chi tiết khu vực trồng trọt
		\item \texttt{PUT /zones/\{zone\_id\}} -- cập nhật khu vực trồng trọt
	\end{itemize}
	
	\noindent
	\textbf{Cảm biến và dữ liệu đo}
	\begin{itemize}
		\item \texttt{GET /sensors?zone\_id=\{id\}} -- danh sách cảm biến theo khu vực
		\item \texttt{POST /sensors} -- thêm cảm biến
		\item \texttt{PATCH /sensors/\{sensor\_id\}} -- cập nhật cảm biến
		\item \texttt{DELETE /sensors/\{sensor\_id\}} -- xóa cảm biến
		\item \texttt{GET /readings} -- lấy dữ liệu đo theo cảm biến và khoảng thời gian
		\item \texttt{GET /readings/latest?zone\_id=\{id\}} -- lấy dữ liệu mới nhất của khu vực
	\end{itemize}
	
	\noindent
	\textbf{Thiết bị điều khiển và thiết bị thông minh}
	\begin{itemize}
		\item \texttt{GET /actuators?zone\_id=\{id\}} -- danh sách thiết bị điều chỉnh
		\item \texttt{POST /actuators} -- thêm thiết bị điều chỉnh
		\item \texttt{PATCH /actuators/\{actuator\_id\}} -- cập nhật thiết bị điều chỉnh
		\item \texttt{DELETE /actuators/\{actuator\_id\}} -- xóa thiết bị điều chỉnh
		\item \texttt{POST /actuators/\{actuator\_id\}/command} -- gửi lệnh điều khiển
		\item \texttt{GET /actuators/\{actuator\_id\}/commands} -- lịch sử lệnh điều khiển
		\item \texttt{POST /devices} -- thêm thiết bị thông minh
		\item \texttt{GET /devices} -- danh sách thiết bị thông minh
		\item \texttt{GET /devices/\{device\_id\}} -- chi tiết thiết bị thông minh
		\item \texttt{PUT /devices/\{device\_id\}} -- cập nhật thiết bị thông minh
		\item \texttt{DELETE /devices/\{device\_id\}} -- xóa thiết bị thông minh
	\end{itemize}
	
	\noindent
	\textbf{Cảnh báo và tổng quan}
	\begin{itemize}
		\item \texttt{GET /alerts} -- danh sách cảnh báo
		\item \texttt{GET /alerts/summary} -- thống kê cảnh báo
		\item \texttt{PATCH /alerts/\{alert\_id\}/read} -- đánh dấu cảnh báo đã đọc
		\item \texttt{GET /dashboard/summary} -- tổng quan khu vực trồng trọt và trạng thái cảm biến
	\end{itemize}
	
	\noindent
	\textbf{Quản trị}
	\begin{itemize}
		\item \texttt{GET /admin/users} -- danh sách nông dân
		\item \texttt{POST /admin/users} -- tạo tài khoản nông dân
		\item \texttt{PATCH /admin/users/\{user\_id\}/status} -- khóa hoặc mở khóa tài khoản
		\item \texttt{GET /admin/zones} -- danh sách toàn bộ khu vực
		\item \texttt{POST /admin/zones} -- tạo khu vực trồng trọt
		\item \texttt{PUT /admin/zones/\{zone\_id\}} -- cập nhật khu vực
		\item \texttt{DELETE /admin/zones/\{zone\_id\}} -- xóa khu vực
		\item \texttt{GET /admin/zones/\{zone\_id\}/farmers} -- danh sách nông dân được phân công
		\item \texttt{POST /admin/zones/\{zone\_id\}/farmers} -- phân công nông dân vào khu vực
		\item \texttt{DELETE /admin/zones/\{zone\_id\}/farmers/\{farmer\_id\}} -- gỡ phân công nông dân
	\end{itemize}
	
	\phantomsection
	\addcontentsline{toc}{subsection}{Phụ lục B: Script cơ sở dữ liệu (PostgreSQL)}
	\subsection*{Phụ lục B: Script cơ sở dữ liệu (PostgreSQL)}
	
	\subsubsection*{B.1. Tạo database}
	
	\noindent
	\begin{lstlisting}[language=SQL]
		CREATE DATABASE agri_db
		WITH ENCODING 'UTF8'
		TEMPLATE template0;
	\end{lstlisting}
	
	\subsubsection*{B.2. Các bảng chính}
	
	\noindent
	Hệ thống gồm các bảng dữ liệu chính sau:
	
	\begin{itemize}
		\item \texttt{users} -- người dùng hệ thống
		\item \texttt{plant\_profiles} -- hồ sơ cây trồng và ngưỡng môi trường
		\item \texttt{growing\_zones} -- khu vực trồng trọt
		\item \texttt{farmer\_zone\_assignments} -- phân công nông dân vào khu vực
		\item \texttt{sensors} -- cảm biến
		\item \texttt{sensor\_readings} -- dữ liệu đo từ cảm biến
		\item \texttt{actuators} -- thiết bị điều chỉnh
		\item \texttt{actuator\_commands} -- lịch sử lệnh điều khiển
		\item \texttt{devices} -- thiết bị thông minh
		\item \texttt{alerts} -- cảnh báo môi trường
	\end{itemize}
	
	Các bảng được thiết kế theo mô hình quan hệ, trong đó \texttt{users} liên kết với \texttt{farmer\_zone\_assignments}; \texttt{growing\_zones} liên kết với \texttt{plant\_profiles}, \texttt{sensors}, \texttt{actuators}, \texttt{devices} và \texttt{alerts}; dữ liệu đo được lưu tại \texttt{sensor\_readings} để phục vụ thống kê, cảnh báo và điều khiển tự động.
	
	\subsubsection*{B.3. Script minh họa một số bảng}
	
	\noindent
	\begin{lstlisting}[language=SQL]
		CREATE TABLE users (
		id UUID PRIMARY KEY,
		email VARCHAR(255) UNIQUE NOT NULL,
		hashed_password VARCHAR(255) NOT NULL,
		full_name VARCHAR(255) NOT NULL,
		role VARCHAR(20) NOT NULL DEFAULT 'user',
		status VARCHAR(20) NOT NULL DEFAULT 'active',
		is_verified BOOLEAN NOT NULL DEFAULT FALSE,
		created_at TIMESTAMP WITH TIME ZONE,
		updated_at TIMESTAMP WITH TIME ZONE
		);
		
		CREATE TABLE plant_profiles (
		id UUID PRIMARY KEY,
		name VARCHAR(255) NOT NULL,
		description TEXT,
		temp_min DOUBLE PRECISION,
		temp_max DOUBLE PRECISION,
		humidity_min DOUBLE PRECISION,
		humidity_max DOUBLE PRECISION,
		soil_moisture_min DOUBLE PRECISION,
		soil_moisture_max DOUBLE PRECISION,
		light_min DOUBLE PRECISION,
		light_max DOUBLE PRECISION,
		ph_min DOUBLE PRECISION,
		ph_max DOUBLE PRECISION,
		ec_min DOUBLE PRECISION,
		ec_max DOUBLE PRECISION,
		growth_period_days INTEGER
		);
		
		CREATE TABLE growing_zones (
		id UUID PRIMARY KEY,
		name VARCHAR(255) NOT NULL,
		description TEXT,
		location VARCHAR(255),
		area_sqm DOUBLE PRECISION,
		plant_profile_id UUID REFERENCES plant_profiles(id),
		is_active BOOLEAN NOT NULL DEFAULT TRUE,
		planting_date DATE,
		expected_harvest_date DATE
		);
		
		CREATE TABLE sensors (
		id UUID PRIMARY KEY,
		name VARCHAR(255) NOT NULL,
		sensor_type VARCHAR(50) NOT NULL,
		unit VARCHAR(50),
		location VARCHAR(255),
		device_address VARCHAR(255),
		update_interval_seconds INTEGER,
		zone_id UUID NOT NULL REFERENCES growing_zones(id),
		is_active BOOLEAN NOT NULL DEFAULT TRUE
		);
		
		CREATE TABLE sensor_readings (
		id UUID PRIMARY KEY,
		sensor_id UUID NOT NULL REFERENCES sensors(id),
		value DOUBLE PRECISION NOT NULL,
		recorded_at TIMESTAMP WITH TIME ZONE NOT NULL
		);
		
		CREATE TABLE alerts (
		id UUID PRIMARY KEY,
		zone_id UUID REFERENCES growing_zones(id),
		sensor_id UUID REFERENCES sensors(id),
		alert_type VARCHAR(50),
		severity VARCHAR(50),
		parameter VARCHAR(50),
		actual_value DOUBLE PRECISION,
		threshold_value DOUBLE PRECISION,
		message TEXT,
		recommended_action TEXT,
		is_read BOOLEAN NOT NULL DEFAULT FALSE,
		triggered_at TIMESTAMP WITH TIME ZONE
		);
	\end{lstlisting}
	
	\phantomsection
	\addcontentsline{toc}{subsection}{Phụ lục C: Code minh họa Backend (FastAPI)}
	\subsection*{Phụ lục C: Code minh họa Backend (FastAPI)}
	
	\subsubsection*{C.1. Controller}
	
	\noindent
	Controller tiếp nhận request REST, kiểm tra dữ liệu đầu vào, lấy thông tin người dùng hiện tại và gọi service để xử lý nghiệp vụ.
	
	\noindent
	\begin{lstlisting}[language=Python]
		@router.get(
		"/summary",
		response_model=BaseResponse[DashboardSummary],
		summary="Tong quan khu vuc trong trot va trang thai cam bien",
		)
		async def get_dashboard_summary(
		current_user: CurrentUser,
		farmer_id: UUID | None = None,
		db: AsyncSession = Depends(get_db),
		):
		data = await service.get_dashboard_summary(current_user, farmer_id, db)
		return BaseResponse.ok(data=data)
	\end{lstlisting}
	
	\subsubsection*{C.2. Service}
	
	\noindent
	Service chứa logic nghiệp vụ chính như tạo khu vực, phân công nông dân, ghi nhận dữ liệu cảm biến, sinh cảnh báo và kích hoạt thiết bị khi thông số môi trường vượt ngưỡng.
	
	\noindent
	\begin{lstlisting}[language=Python]
		async def list_zones(self, user: User) -> list[GrowingZone]:
		query = select(GrowingZone)
		if user.role != UserRole.ADMIN:
		query = query.join(FarmerZoneAssignment).where(
		FarmerZoneAssignment.farmer_id == user.id
		)
		result = await self.db.execute(query)
		return list(result.scalars().all())
	\end{lstlisting}
	
	\subsubsection*{C.3. Security}
	
	\noindent
	Backend sử dụng JWT, bcrypt và dependency của FastAPI để xác thực, phân quyền theo vai trò và bảo vệ các endpoint cần đăng nhập.
	
	\noindent
	\begin{lstlisting}[language=Python]
		pwd_context = CryptContext(schemes=["bcrypt"], deprecated="auto")
		
		def hash_password(password: str) -> str:
		return pwd_context.hash(password)
		
		def create_access_token(subject: Any, extra: dict = None) -> str:
		payload = {
			"sub": str(subject),
			"type": TokenType.ACCESS,
			"exp": expire,
			"iat": datetime.now(timezone.utc),
		}
		return jwt.encode(payload, settings.JWT_SECRET_KEY, algorithm=settings.JWT_ALGORITHM)
	\end{lstlisting}
	
	\subsubsection*{C.4. Response chuẩn}
	
	\noindent
	Lớp \texttt{BaseResponse\textless T\textgreater} giúp chuẩn hóa kết quả trả về, bảo đảm frontend luôn nhận cùng một cấu trúc dữ liệu.
	
	\noindent
	\begin{lstlisting}[language=Python]
		class BaseResponse(BaseModel, Generic[T]):
		success: bool = True
		message: str = "Thanh cong"
		data: Optional[T] = None
		
		@classmethod
		def ok(cls, data: T = None, message: str = "Thanh cong"):
		return cls(success=True, message=message, data=data)
	\end{lstlisting}
	
	\phantomsection
	\addcontentsline{toc}{subsection}{Phụ lục D: Frontend (React + TypeScript)}
	\subsection*{Phụ lục D: Frontend (React + TypeScript)}
	
	\subsubsection*{D.1. Cấu trúc dự án}
	
	\noindent
	\begin{itemize}
		\item \texttt{components} -- thành phần giao diện dùng chung
		\item \texttt{routes} -- định nghĩa màn hình và điều hướng
		\item \texttt{features} -- nhóm chức năng theo nghiệp vụ
		\item \texttt{hooks} -- hook gọi API và xử lý dữ liệu
		\item \texttt{stores} -- quản lý trạng thái người dùng và giao diện
		\item \texttt{services} -- cấu hình Axios và endpoint API
		\item \texttt{types} -- kiểu dữ liệu TypeScript
		\item \texttt{lib} -- hằng số, mapper và schema kiểm tra dữ liệu
	\end{itemize}
	
	\subsubsection*{D.2. Đặc điểm chính}
	
	\noindent
	\begin{itemize}
		\item React kết hợp TypeScript để kiểm soát kiểu dữ liệu khi phát triển giao diện.
		\item TanStack Router tổ chức route cho các màn hình đăng nhập, tổng quan, khu vực, cảm biến, thiết bị, cảnh báo và người dùng.
		\item Axios được cấu hình \texttt{baseURL} là \texttt{/api/v1}, timeout 15 giây và \texttt{withCredentials} để gửi cookie xác thực.
		\item Zustand quản lý trạng thái đăng nhập và trạng thái giao diện.
		\item Các hook như \texttt{useZones}, \texttt{useSensors}, \texttt{useDevices}, \texttt{useAlerts} đóng vai trò trung gian giữa giao diện và API.
	\end{itemize}
	
	\subsubsection*{D.3. Cấu hình Axios}
	
	\noindent
	\begin{lstlisting}
		export const api = axios.create({
			baseURL: '/api/v1',
			timeout: 15000,
			headers: {
				'Content-Type': 'application/json',
			},
			withCredentials: true,
		})
		
		api.interceptors.response.use(
		(response) => response,
		async (error) => {
			if (error.response?.status === 401) {
				useAuthStore.getState().logout()
				window.location.href = '/login'
			}
			return Promise.reject(error)
		},
		)
	\end{lstlisting}
	
	\phantomsection
	\addcontentsline{toc}{subsection}{Phụ lục E: Bảo mật hệ thống}
	\subsection*{Phụ lục E: Bảo mật hệ thống}
	
	\subsubsection*{E.1. Cơ chế xác thực}
	
	\noindent
	\begin{itemize}
		\item Người dùng đăng nhập bằng email và mật khẩu.
		\item Backend tạo JWT access token sau khi xác thực thành công.
		\item Token được lưu trong HTTP-only cookie nhằm hạn chế truy cập trực tiếp từ JavaScript phía client.
		\item Các API nghiệp vụ yêu cầu người dùng hợp lệ thông qua dependency \texttt{CurrentUser}.
	\end{itemize}
	
	\subsubsection*{E.2. Phân quyền}
	
	\noindent
	\begin{itemize}
		\item \texttt{ADMIN}: quản lý người dùng, hồ sơ cây trồng, khu vực trồng trọt và phân công nông dân.
		\item \texttt{USER}: theo dõi khu vực được phân công, quản lý cảm biến, thiết bị và cảnh báo trong phạm vi được cấp quyền.
	\end{itemize}
	
	\subsubsection*{E.3. Bảo mật hệ thống}
	
	\noindent
	\begin{itemize}
		\item Mật khẩu được mã hóa bằng bcrypt trước khi lưu vào cơ sở dữ liệu.
		\item JWT ký bằng \texttt{JWT\_SECRET\_KEY} và thuật toán \texttt{HS256}.
		\item CORS giới hạn domain được phép truy cập thông qua biến cấu hình \texttt{ALLOWED\_ORIGINS}.
		\item FastAPI và Pydantic kiểm tra dữ liệu đầu vào theo schema.
		\item SQLAlchemy xây dựng truy vấn có tham số, hạn chế nguy cơ SQL Injection.
	\end{itemize}
	
	\subsubsection*{E.4. Chống tấn công}
	
	\noindent
	\begin{itemize}
		\item SQL Injection: sử dụng ORM SQLAlchemy và truy vấn tham số hóa.
		\item XSS: hạn chế render HTML thô ở frontend, kiểm tra dữ liệu nhập bằng schema.
		\item CSRF: cookie xác thực đặt \texttt{SameSite=Lax}; môi trường production cần bật HTTPS và cấu hình \texttt{secure}.
		\item Truy cập trái phép: kiểm tra vai trò \texttt{ADMIN} hoặc \texttt{USER} trước khi thực hiện thao tác nhạy cảm.
	\end{itemize}
	
	\phantomsection
	\addcontentsline{toc}{subsection}{Phụ lục F: Triển khai hệ thống}
	\subsection*{Phụ lục F: Triển khai hệ thống}
	
	\subsubsection*{F.1. Kiến trúc hệ thống}
	
	\noindent
	\begin{itemize}
		\item React frontend
		\item FastAPI backend
		\item PostgreSQL database
		\item Redis cache
		\item MailHog phục vụ kiểm thử email trong môi trường phát triển
		\item Docker Compose quản lý các container
	\end{itemize}
	
	\subsubsection*{F.2. Các service trong Docker Compose}
	
	\noindent
	\begin{itemize}
		\item \texttt{server}: backend FastAPI, chạy tại cổng \texttt{8000}
		\item \texttt{fe}: frontend React, chạy tại cổng \texttt{5176}
		\item \texttt{db}: PostgreSQL 16, ánh xạ cổng \texttt{5436}
		\item \texttt{redis}: Redis 7, ánh xạ cổng \texttt{6379}
		\item \texttt{mailhog}: kiểm thử email, giao diện tại cổng \texttt{8025}
	\end{itemize}
	
	\subsubsection*{F.3. Quy trình triển khai}
	
	\noindent
	\begin{lstlisting}
		docker compose up --build
		docker compose logs -f server
		docker compose down
	\end{lstlisting}
	
	\noindent
	Sau khi các container khởi động, người dùng truy cập frontend qua \texttt{http://localhost:5176}, backend qua \texttt{http://localhost:8000}, tài liệu API ở \texttt{http://localhost:8000/docs} trong môi trường phát triển.
	
	\subsubsection*{F.4. Giám sát và bảo trì}
	
	\noindent
	\begin{itemize}
		\item Kiểm tra log backend bằng \texttt{docker compose logs -f server}.
		\item Sao lưu cơ sở dữ liệu PostgreSQL định kỳ bằng \texttt{pg\_dump}.
		\item Theo dõi trạng thái container bằng \texttt{docker compose ps}.
		\item Quản lý biến môi trường qua tệp \texttt{be/.env}.
	\end{itemize}
	
	\noindent
	\textbf{Kết luận phụ lục}
	
	Phụ lục cung cấp tài liệu kỹ thuật tổng hợp của hệ thống, bao gồm API, cơ sở dữ liệu, backend, frontend, bảo mật và triển khai. Các thành phần được tổ chức theo kiến trúc tách lớp, giúp hệ thống dễ bảo trì, dễ mở rộng và đảm bảo tính nhất quán giữa giao diện người dùng, xử lý nghiệp vụ và lưu trữ dữ liệu.
	
\end{document}
